% Options for packages loaded elsewhere
\PassOptionsToPackage{unicode}{hyperref}
\PassOptionsToPackage{hyphens}{url}
%
\documentclass[
  11pt,
]{article}
\usepackage{amsmath,amssymb}
\usepackage{iftex}
\ifPDFTeX
  \usepackage[T1]{fontenc}
  \usepackage[utf8]{inputenc}
  \usepackage{textcomp} % provide euro and other symbols
\else % if luatex or xetex
  \usepackage{unicode-math} % this also loads fontspec
  \defaultfontfeatures{Scale=MatchLowercase}
  \defaultfontfeatures[\rmfamily]{Ligatures=TeX,Scale=1}
\fi
\usepackage{lmodern}
\ifPDFTeX\else
  % xetex/luatex font selection
\fi
% Use upquote if available, for straight quotes in verbatim environments
\IfFileExists{upquote.sty}{\usepackage{upquote}}{}
\IfFileExists{microtype.sty}{% use microtype if available
  \usepackage[]{microtype}
  \UseMicrotypeSet[protrusion]{basicmath} % disable protrusion for tt fonts
}{}
\makeatletter
\@ifundefined{KOMAClassName}{% if non-KOMA class
  \IfFileExists{parskip.sty}{%
    \usepackage{parskip}
  }{% else
    \setlength{\parindent}{0pt}
    \setlength{\parskip}{6pt plus 2pt minus 1pt}}
}{% if KOMA class
  \KOMAoptions{parskip=half}}
\makeatother
\usepackage{xcolor}
\usepackage[margin=1in]{geometry}
\usepackage{longtable,booktabs,array}
\usepackage{calc} % for calculating minipage widths
% Correct order of tables after \paragraph or \subparagraph
\usepackage{etoolbox}
\makeatletter
\patchcmd\longtable{\par}{\if@noskipsec\mbox{}\fi\par}{}{}
\makeatother
% Allow footnotes in longtable head/foot
\IfFileExists{footnotehyper.sty}{\usepackage{footnotehyper}}{\usepackage{footnote}}
\makesavenoteenv{longtable}
\setlength{\emergencystretch}{3em} % prevent overfull lines
\providecommand{\tightlist}{%
  \setlength{\itemsep}{0pt}\setlength{\parskip}{0pt}}
\setcounter{secnumdepth}{-\maxdimen} % remove section numbering
\usepackage{amsmath,amssymb}
\ifLuaTeX
  \usepackage{selnolig}  % disable illegal ligatures
\fi
\IfFileExists{bookmark.sty}{\usepackage{bookmark}}{\usepackage{hyperref}}
\IfFileExists{xurl.sty}{\usepackage{xurl}}{} % add URL line breaks if available
\urlstyle{same}
\hypersetup{
  pdftitle={The Value of Stopping: Evidence Queues, Recoverability, and Priority Rights},
  pdfauthor={Dr.~Jan-Hendrik Weinert},
  hidelinks,
  pdfcreator={LaTeX via pandoc}}

\title{The Value of Stopping: Evidence Queues, Recoverability, and
Priority Rights}
\author{Dr. Jan-Hendrik Weinert}
\date{19 October 2026}


% Publication formatting edition, 19 October 2026. Scientific equations unchanged.
\usepackage{xurl,needspace,etoolbox}
\geometry{a4paper,margin=25.4mm}
\hypersetup{hidelinks,pdfauthor={Dr. Jan-Hendrik Weinert}}
\setlength{\parindent}{1.2em}
\setlength{\parskip}{3pt plus 1pt}
\setlength{\emergencystretch}{3em}
\hfuzz=0.2pt
\raggedbottom
\pagestyle{plain}
\makeatletter
\renewcommand{\@seccntformat}[1]{\csname the#1\endcsname\ifstrequal{#1}{section}{.}{}\quad}
\renewcommand{\maketitle}{%
  \begin{center}
  {\LARGE\bfseries\@title\par}
  \vspace{0.8em}
  {\normalsize\@author\par}
  \vspace{0.35em}
  {\small\@date\par}
  \end{center}
  \vspace{0.4em}}
\renewenvironment{abstract}{%
  \begin{center}\small\bfseries Abstract\end{center}
  \begin{list}{}{\setlength{\leftmargin}{0.5in}\setlength{\rightmargin}{0.5in}}
  \item[]\small\setlength{\parindent}{0pt}\ignorespaces
}{\end{list}\vspace{0.6em}}
\makeatother

\begin{document}
\maketitle
\begin{abstract}
Correlated evidence can make efficient research privately unattractive both initially and at continuation. In a two-case, two-slot model with expiration and optional terminal action, we derive the exact implementation and funding frontier for two nonnegative public-history credits. One rewards the initial experiment; the other rewards a specified contingent final experiment. Complete backward induction gives necessary and sufficient inequalities, separate expected and pathwise gross-funding limits, and attainment boundaries on a certified interval of binary priors. Credits are transfers rather than resource losses. The independent benchmark retains the recoverability-versus-priority comparison and common-charge classification. Complete proofs and exact synthetic enumeration support this restricted implementation result; no unrestricted mechanism optimum or empirical effect is claimed.
\end{abstract}
\noindent\small \textbf{Keywords:} evidence acquisition; scarce attention; decision
queues; stopping; priority rights; expiration.\par\normalsize
\vspace{0.6em}
\section{1. Introduction}\label{introduction}

A request for one more analysis can be justified for its own decision
and still be costly for the organization. The analyst may have to
postpone research on another case whose purchasing quote expires sooner,
whose decision remains unusually uncertain, or whose recoverable value
will be lost if action is delayed. The relevant opportunity cost is
therefore endogenous to the competing case. It is not generally a fixed
price for elapsed time.

We study two evidence cases sharing one analyst for two research slots.
Each case can be acted on immediately without new evidence, or receive
one pre-specified experiment before action. While a case is kept
unresolved, holding costs accrue and its opportunity may expire. A case
assigned the second slot can instead be closed immediately. Thus the
allocation of attention and the decision to stop researching must be
determined jointly. The experiment is a bounded evidence packet, not an
open-ended search process. This finite restriction makes the
organizational equilibrium, corrective rule, and implementation
boundaries fully solvable.

For case \(i\), let \(b_i\) be the net gain from immediate research
rather than immediate action. Let \(w_i\) be the optimal gain from being
offered research only after one waiting slot, allowing immediate action
instead. Then \(u_i=b_i^+-w_i\) is the loss from delaying that research
option. The efficient order compares \(u_i\), not just \(b_i\). This
loss depends on beliefs, signal precision, stakes, recoverability,
holding costs, and expiration.

We analyze a concrete decentralized institution: case A initially owns
the right to the first slot. Its owner can claim that slot, release it
and wait, or act without research. Case B then receives the available
slot; final-slot choices are privately optimal. Owners value their own
case's payoff and pay their own research and holding costs. Even with
complete information and no split accounting of these direct costs, A
generally claims early service whenever its own research has positive
net value. It ignores how much B loses by waiting.

The independent benchmark contributes a
recoverability-versus-deferral-priority comparison and the exact
restrictions under which a common research charge can implement
efficient retention and order. The underlying scheduling interchange and
externality toll are classical. We give a proposition-by-proposition
comparison rather than claiming a new Gittins index or dynamic
mechanism. We cannot certify that the bounded increments are unoccupied
because one close predecessor remains unavailable in full text.

We therefore also solve the smallest adjacent nonseparable margin, still
with two cases, one experiment each, and two slots. An explicitly
correlated prior makes the first public signal change the other case's
research value. A previously researched case may remain unresolved to
benefit from the second signal. This creates a cross-information benefit
as well as congestion. We derive all contingent joint choices, the
corresponding finite local game, and the exact final-slot boundary at
which an information credit rather than a positive toll is needed. The
central correlated result then closes the initial incentive gap using
just two observable-history credits. We solve the complete local game,
give exact necessary and sufficient implementation inequalities, and
distinguish expected funding from pathwise capacity, including equality
and nonattainment. An exact interval of priors around the original
witness makes this a parameter-region result. Neither correlation nor
conditional subsidies are claimed as new principles.

\section{2. Literature and contribution
boundary}\label{literature-and-contribution-boundary}

Henry and Ottaviani (2019) study costly sequential research under
divided control and conflicting terminal interests. Escobar and Zhang
(2021) analyze private learning, communication, and prompt
implementation in dynamic delegation. They already establish that
organizational rules can alter learning and delay. We instead align
owners with their own case's economic payoff, make primitives public,
and introduce competition between cases for a single analyst. Our
friction does not require persuasion, private signals withheld from the
organization, or different action preferences within a case.

Sequential sampling is the separable foundation, not a novelty claim.
Fudenberg, Strack, and Strzalecki (2018) connect costly information
gathering to decision timing. Xiao, Hu, and Wang (2026), published
online in \emph{Management Science}, characterize recursive standards in
delegated finite search. Their publisher abstract and publication status
were verified; full technical access remained unavailable. We cannot
certify priority against every result in that paper.

The relevant scheduling comparison is broader than sequential sampling.
Gittins (1979) develops dynamic allocation indices, with applications to
scheduling and search; Weber (1992) provides a proof of the classical
index policy. Those results are much more general than our two-job
calculation in important respects. Our \(u_i\) is a finite pairwise
interchange score, not a renamed Gittins benefit/time ratio. Inactive
evidence cases can expire and the common finite horizon shrinks, so a
general rested-bandit theorem cannot simply be invoked to justify an
index policy for extensions of our model.

Expiring-arm and dynamic-pivot precedents are mapped explicitly in
Sections 2.1-2.3. Whittle (1988), Weber and Weiss (1990), Righter and
Shanthikumar (1998), and Yu, Xu, and Tong (2017 preprint; later 2018
journal publication) are scheduling comparators. Adiri and Yechiali
(1974) and Bergemann and Välimäki (2010) are pricing precedents. The
posted public-state toll is not claimed to duplicate a complete
private-information pivot mechanism.

Campbell, Ederer, and Spinnewijn (2014) study effort and communication
around deadlines. Their mechanism cautions against interpreting our
no-effort-response deadline result as universal. Eyring, Ferguson, and
Koppers (2021) vary the composition of performance feedback, not
acquisition or queue priority. Their experiment supports separating
display from acquisition interventions. Xu, Zariphopoulou, and Zhang
(2026, arXiv v4) study reversible choices with changing information
regimes. Our recoverability parameter is more limited: it modifies the
wrong-action payoff without creating post-action learning.

The primary source ledger records the depth of reading and version
differences. Full text was read for several close organizational and
pricing papers and the deadline-scheduling preprint; original Gittins,
Righter-Shanthikumar, and recent delegated-search technical coverage
remains incomplete. We make no exhaustive-literature or
established-priority claim.

\subsection{2.1 Primitive mapping to expiring-arm
scheduling}\label{primitive-mapping-to-expiring-arm-scheduling}

Our independent-case state is \((p_i,e_i,l_i,t)\): belief, remaining
experiment flag, live/retired status, and the common clock. An active
arm pays \(-c_i-d_i\), faces expiry, and if surviving generates its one
signal. A passive retained arm pays \(-d_i\), faces expiry, and keeps
its belief. Retirement pays \(g_i(p_i)\) once and removes the case;
expiry pays zero. Thus this is a finite-horizon expiring-arm stopping
problem. Neither expiry nor the optional stopping action is a new
scheduling principle.

Subtracting the potential \(g_i\) makes the correspondence precise. The
one-slot expected transformed reward is \(b_i\) when active, \(-\ell_i\)
when passive, and zero when retired. The initial potential is added back
once. This is just accounting; it does not identify a new index or prove
indexability. The clock and passive expiry remain state changes even
when the belief is frozen. The optional immediate payoff is not a reason
to declare the problem outside every existing bandit framework: it can
be encoded as a terminal action. Time augmentation likewise accommodates
a finite horizon, although it need not preserve a rested-arm index
theorem's assumptions.

Righter-Shanthikumar's publisher abstract confirms index results under
independent expiry. A lawful full version was not retrieved after
checking the publisher and both authors' public bibliographies. We did
\textbf{not} inspect their exact active/passive hazard assumptions, arm
transition/reward restrictions, horizon/discount criterion, retirement
options, information structure, or index optimality conditions on pages
453-468. None of those cells is marked verified below. We therefore make
no claim that their theorem excludes or includes our model. An abstract
cannot supply that comparison.

Whittle's publisher description concerns relaxed capacity prices and
active/passive evolution. The author-hosted Weber-Weiss full text
establishes that asymptotic index conclusions require additional
conditions. The read Yu-Xu-Tong v5 preprint treats stochastic deadlines,
capacity, costs and noncompletion penalties; it is not the same read
object as its later journal publication. These precedents establish
occupied problem classes, not a theorem-by-theorem novelty gap for this
paper.

\subsection{2.2 Primitive mapping to dynamic-pivot
pricing}\label{primitive-mapping-to-dynamic-pivot-pricing}

Bergemann-Välimäki's published Section 3 allocates a central facility,
one task per period, with delay represented by discounting. Section 4,
equations (5)-(7), prices current and continuation externalities using
programs with and without an agent. Their independent private-value
Markov framework yields truth-telling and participation results that are
not proved here.

For our two independent routes, B's optimal incremental value if A
releases first service is \(b_B^+\); if A occupies it, B retains or acts
and obtains \(w_B\). Their difference \(u_B\) is exactly the marginal
externality algebra, with the common baseline \(g_B\) canceling. There
is no extra loss to add after subtracting those two values. Our
asymmetric right-holder charge and uncharged B are not the complete
dynamic-pivot transfer vector. Internal charges are excluded from firm
resource payoff. A finite absorbing horizon is not, by itself, an
economic novelty relative to dynamic pricing; finite flows can be
represented separately from a discount convention.

In the correlated extension, one case's signal changes the other's
posterior. The independently evolving private-value conditions in that
published theorem are not directly satisfied. This is a boundary of that
direct theorem mapping, not a claim that correlated-information
externality pricing is absent from the wider literature. The extension
keeps information public and makes no private-type truthfulness claim.

The restricted credits in Sections 8.4-8.5 do not replicate realized
team payoffs or the full dynamic-pivot payment vector. The published
pivot equations compute each agent's flow and continuation externality
from with/without-agent programs; our two instruments are constrained to
specified history/action events. Their independent private-value
truthfulness theorem is not being extended to correlated private types.
Known public beliefs, observable research and enforceable credit
commitment suffice here; eventual action-payoff observability is
unnecessary. Their positive payment variable is subtracted from the
agent's utility; our credits enter with the opposite sign. Their
published no-deficit conclusion therefore does not supply our gross
credit budget. The funding frontiers describe internal appropriation
authority, not budget balance or the pivot mechanism's revenue
properties.

Moroni (2022), published introduction and model, studies
history-dependent incentives and free-riding in a multistage risky
project with privately observed effort and learning, limited liability
and endogenous assignments. That precedent further rules out novelty
claims for conditional experimentation rewards. Our observable
one-experiment actions, two scarce slots and two fixed payment events
define a different restricted calculation; we do not infer an
unrestricted contracting advantage from that difference.

\subsection{2.3 Result-by-result
comparison}\label{result-by-result-comparison}

\begin{longtable}[]{@{}
  >{\raggedright\arraybackslash}p{(\columnwidth - 4\tabcolsep) * \real{0.3333}}
  >{\raggedright\arraybackslash}p{(\columnwidth - 4\tabcolsep) * \real{0.3333}}
  >{\raggedright\arraybackslash}p{(\columnwidth - 4\tabcolsep) * \real{0.3333}}@{}}
\toprule\noalign{}
\begin{minipage}[b]{\linewidth}\raggedright
Result in this paper
\end{minipage} & \begin{minipage}[b]{\linewidth}\raggedright
Closest established principle
\end{minipage} & \begin{minipage}[b]{\linewidth}\raggedright
Exact bounded claim and remaining limit
\end{minipage} \\
\midrule\noalign{}
\endhead
\bottomrule\noalign{}
\endlastfoot
Separable benchmark, Section 3 & Standard sequential stopping and
benefit/time accounting & No novelty claimed \\
Theorem 1 & Finite two-job interchange; expiring-arm scheduling &
Explicit optional retirement/retention solution; not a new index;
Righter-Shanthikumar conditions uninspected \\
Theorem 2 and exact fee set & Local queue priority and restricted
pricing & Necessary/sufficient common-fee intervals within the stated
institution and tie convention; no claim all price mechanisms fail \\
Corollary 1 & Enumeration of feasible deadlines & Exact three-deadline
comparison only \\
Theorem 3 and Corollary 2 & Marginal externality and priority pricing &
Public-state asymmetric implementation; not a new dynamic-pivot
theorem \\
Theorem 4 & Bayes-risk payoff comparison combined with
deteriorating-option scheduling & Opposite recoverability derivatives on
the retained-research region; priority crossing is conditional and
priority remains uncertified \\
Theorems 5-6 & Finite Bayesian dynamic programming; information
externalities & Correlated joint solution and final-slot boundary,
retained benchmarks \\
Theorems 7-8 & Conditional subsidies, free-riding compensation,
restricted dynamic transfers & Complete two-history-credit
implementation and piecewise expected/pathwise funding frontier on I; no
unrestricted optimal-contract or pivot claim \\
\end{longtable}

Klein and Rady (2011) already study strategic experimentation with
correlated risky types. The June 23, 2010 author-hosted manuscript and
2011 publisher metadata were checked; their continuous-time two-player
experimentation setting is distinct from our two-slot retirement
problem. Correlation and cross-learning themselves are not claimed new.
The inspected working model permits simultaneous safe/risky choices with
publicly observed payoffs and continuous-time learning, rather than our
single analyst, two optional one-shot experiments and expiring action
opportunities. Its public cross-learning and free-riding logic is
inherited. Theorems 7-8 instead restrict the payment support to two
existing observable research histories, preserve local stopping rights,
and solve their funding limits. No located theorem was verified to state
this exact restricted frontier; that bounded reading is not an
exhaustive priority certification. The uninspected independent-expiry
article is not a positive premise of the correlated-credit result.

\section{3. Separable stopping
benchmark}\label{separable-stopping-benchmark}

To locate the new mechanism, consider first a single finite evidence
tree. At live history \(x\), immediate action pays
\(g(x)=\max_a E[u(a,\theta)\mid\mathcal F_x]\). An experiment costs
\(c_e\), consumes a period with additive delay price \(d\), and survives
with probability \(s_e\). The filtration records observed signals,
research choices, elapsed slots, and opportunity survival. Feasible
policies are adapted and terminate at a finite horizon. Expiration pays
zero. The value obeys

\[V_d(x)=\max\{g(x),\max_e[-c_e-d+s_eE_eV_d(x')]\}.\]

For any adaptive policy \(\pi\) that continues first, define expected
terminal reward \(R_\pi\), acquisition cost \(C_\pi\), and active
periods \(T_\pi>0\). Then continuing is strictly better than acting now
if and only if

\[d<D(x),\qquad D(x)=\max_\pi\frac{R_\pi-C_\pi-g(x)}{T_\pi}.\]

This follows by rewriting each policy's incremental payoff as
\(R_\pi-C_\pi-g-dT_\pi\). It is elementary finite optimal stopping. If
an analyst internalizes only \(\alpha d\), its continuation region uses
\(\alpha d<D(x)\). Charging the missing cost aligns objectives only when
this marginal charge enters actual preferences and work cannot escape
the ledger.

Likewise, Bellman action regret telescopes: optimal value minus any
feasible policy's payoff equals the expected sum of optimal action-value
gaps along that policy's path, including expiration. The complete proofs
and earlier exact single-case illustrations are retained in Appendix B.
They are not the revised paper's central contribution.

A bound \(d\in[d_L,d_U]\) provides a stopping certificate only
conditional on the full payoff, signal, expiration, action, and policy
model. It is not robustness to misspecified likelihoods or missed
alternatives. In the shared-resource models below, no exogenous \(d\)
alone represents the forgone value of another case.

\section{4. Two cases sharing one
analyst}\label{two-cases-sharing-one-analyst}

\subsection{4.1 Timing, primitives, and economic
accounting}\label{timing-primitives-and-economic-accounting}

There are two cases \(i\in\{A,B\}\), independent latent binary states
\(\theta_i\), and public priors \(p_i\). Slots are \(t=0,1\), followed
by mandatory terminal action on surviving unresolved cases at \(t=2\).
One analyst can perform at most one experiment in each slot. Each case
has at most one experiment. Its binary signal has accuracy
\(q_i\in[1/2,1]\), conditionally independent of the other case. All
parameters and available actions are public.

At each slot boundary the organization may act on any subset of
unresolved cases, with no analyst time required for that terminal
action. Correct action yields \(S_i>0\); wrong action yields \(r_iS_i\),
\(0\leq r_i\leq1\). A resolved case leaves permanently. A case retained
unresolved over a slot incurs \(d_i\geq0\) and survives the slot with
independent probability \(s_i=1-h_i\), where \(0\leq h_i<1\). If the
analyst researches case \(i\), it also incurs acquisition cost
\(c_i\geq0\). The signal is received only if the case survives; the
researched surviving case can act immediately at the next boundary.
Since it has no experiments left, additional waiting is weakly dominated
by action.

Expiration pays zero. Costs for an occupied slot are paid before
expiration. Thus the same accounting applies whether a case is active or
passively waiting. The additive holding cost is a residual resource or
operating loss, excluding any loss already counted through the terminal
payoff and expiration. A first-slot research on A with B retained costs
\(c_A+d_A+d_B\), not \(c_A\) twice and not a second artificial
queue-loss penalty. Internal corrective transfers are not firm resource
costs.

The filtration contains the public priors, decisions to stop or wait,
analyst assignment, and observed survival and signals. An unrestricted
joint policy may condition its final-slot assignment on all this
information. Idle policies are feasible and included in the independent
enumeration. The theorem will show that no idle policy improves on the
stated solution. There are no new arrivals, strategic suppliers,
correlated information, hidden effort, or endogenous experimental
precision.

\subsection{4.2 Information value and the option to
retire}\label{information-value-and-the-option-to-retire}

Write \(m_i=\min(p_i,1-p_i)\) and let \(\bar m_i\) be expected posterior
Bayes error after the one available signal. Define

\[g_i=S_i[1-(1-r_i)m_i],\quad
M_i=S_i[1-(1-r_i)\bar m_i],\]

\[b_i=s_iM_i-c_i-d_i-g_i,\qquad
\ell_i=d_i+h_ig_i. \tag{1}\]

Here \(g_i\) is immediate-action value, \(M_i\) the expected payoff
after observing the signal conditional on survival, and \(b_i\) the net
gain from researching in the next available slot. The quantity
\(\ell_i\) is the expected loss from keeping the current decision
unresolved for one slot even without acquiring evidence.

If case \(i\) is offered only the second slot, it may act now or wait
and, conditional on survival, choose whether to research. Its
incremental value relative to immediate action is

\[w_i=[s_ib_i-\ell_i]^+,\qquad z^+=\max(z,0). \tag{2}\]

Since \(s_i\leq1\) and \(\ell_i\geq0\), \(0\leq w_i\leq b_i^+\). Define
the option's deferral loss

\[u_i=b_i^+-w_i\geq0. \tag{3}\]

This loss includes optimal stopping. In particular, when deferred
research is no longer worthwhile, \(w_i=0\) and \(u_i=b_i^+\):
displacing that case removes its entire positive research option, but it
can still act now. It is incorrect to add its immediate payoff \(g_i\)
as a second loss after using this expression.

\section{5. Optimal joint stopping and
scheduling}\label{optimal-joint-stopping-and-scheduling}

\textbf{Theorem 1 (independent two-case solution).} Let \(G=g_A+g_B\).
The optimal expected total payoff is

\[W^*=G+\max\{0,\ b_A+w_B,\ b_B+w_A\}. \tag{4}\]

If both \(b_i>0\), researching A first is optimal rather than
researching B first exactly when \(u_A\geq u_B\). After assigning \(i\)
first, the other case \(j\) waits precisely when \(s_jb_j>\ell_j\);
otherwise it acts immediately. A surviving second-slot case is
researched if \(b_j>0\), otherwise acted on. After the single signal,
the researched case selects its posterior-optimal terminal action.

\textbf{Proof.} If case \(i\) is researched first, its expected
incremental payoff is \(b_i\). Independence implies that the other case
learns nothing about its own state from \(i\)'s signal or expiration.
Its best alternative to immediate action is one slot of waiting followed
by its own optimal research-or-act choice. That incremental payoff is
\(w_j\) by (2). Therefore the best value conditional on \(i\) first is
\(G+b_i+w_j\). Stopping both gives \(G\).

An idle first slot leaves at most one later experiment. Each retained
case loses \(\ell_i\) from baseline while waiting. If at most one
retained case survives, its later research gain cannot exceed its
immediate positive research gain. If both survive, the final allocation
selects at most one. Removing the idle slot and performing that chosen
experiment immediately is no worse in expectation: the expected maximum
of the gains of surviving cases is at most \(\max(b_A^+,b_B^+)\), and
waiting losses are nonnegative. Immediate research on a maximizer and
immediate action on the other case attains \(G+\max(b_A^+,b_B^+)\).
Hence idle policies cannot beat the candidates in (4). No other
first-slot policy exists.

When both \(b_i>0\), subtracting the two research-order values gives
\((b_A-w_A)-(b_B-w_B)=u_A-u_B\). The retention and final action rules
follow from the same comparison. \(\square\)

The formula is a finite scheduling interchange, not a general bandit
index theorem. In this model signals affect posterior terminal action
but do not create further research choices on the same case.
Independence therefore eliminates informational cross-effects in the
final schedule. The stochastic tree still includes both signal outcomes
and both cases' expiration; a stronger multi-experiment model would
generally need a joint Bellman solution and is not claimed solved here.

\section{6. Decentralized priority rights and corrective
policies}\label{decentralized-priority-rights-and-corrective-policies}

\subsection{6.1 Local equilibrium}\label{local-equilibrium}

Case A owns the first-slot priority right. Its owner may research now,
act now, or release the slot and wait while B chooses whether to
research. If A claims the slot, B chooses whether to act now or wait for
the second slot. Subsequent research and terminal actions are locally
optimal. Owners maximize their own case's expected payoff; no owner
receives the other's payoff. Internal transfers enter their objectives
linearly. Ties between acting and costly research favor action; ties
between using an available slot now and waiting with identical payoff
favor earlier service. These conventions and public primitives define a
perfect-information sequential game.

A common nonnegative charge \(\tau\) is paid for every experiment,
regardless of when it occurs, to a central budget. Let
\(\tilde b_i=b_i-\tau\) and

\[\tilde w_i=[s_i(b_i-\tau)-\ell_i]^+.\]

\textbf{Theorem 2 (independent local equilibrium and common-charge
boundary).} Under the stated tie convention, if \(b_A>\tau\), A
researches first and B waits if and only if \(\tilde w_B>0\). If
\(b_A\leq\tau\), A acts immediately and B researches if and only if
\(b_B>\tau\). When \(\ell_A>0\), A's preference for immediate research
over waiting is strict whenever \(b_A>\tau\).

Suppose the unique efficient order is B first, \(b_A,b_B>0\), and
\(\ell_A>0\).

\begin{enumerate}
\def\labelenumi{(\alph{enumi})}
\item
  If \(w_A>0\), no common \(\tau\geq0\) implements the efficient
  schedule and retention policy.
\item
  If \(w_A=0\), every \(\tau\in[b_A,b_B)\) implements it; this interval
  is nonempty.
\end{enumerate}

If A first is efficient, \(\tau=0\) implements an efficient policy. If
\(b_A\leq0\), \(\tau=0\) lets B research when worthwhile. Thus failure
is conditional, not universal.

\textbf{Proof.} With positive \(\tilde b_A\), immediate research gives A
gain \(\tilde b_A\). If A releases the first slot, its best waiting gain
when B uses it is \(\tilde w_A\leq\tilde b_A\). The difference before
truncation is \(h_A\tilde b_A+\ell_A\), positive if \(\ell_A>0\). If B
does not use the released slot, idling provides no improvement. Hence A
claims service. If \(\tilde b_A\leq0\), waiting cannot help and A acts.
B's behavior follows (2) with the charged research cost.

For (a), reversing A's strictly preferred first claim requires
\(\tau\geq b_A\). But then \(b_A-\tau\leq0\), so A will not research in
the last slot and cannot have positive retained research value. If
\(\tau<b_A\), A takes the first slot. Neither case gives the desired
policy.

For (b), unique B-first optimality and \(w_A=0\) imply
\(b_B>b_A+w_B\geq b_A\). The displayed interval makes A act and B
research. The other implementation claims follow from the equilibrium at
\(\tau=0\). \(\square\)

For an exact classification, let \(t_B=b_B-\ell_B/s_B\). Excluding
transfers from resource payoff, the equilibrium value at every
\(\tau\geq0\) is

\[J_\tau=\begin{cases}
G+b_A+\mathbf1\{\tau<t_B\}w_B,&\tau<b_A,\\
G+\mathbf1\{\tau<b_B\}b_B,&\tau\geq b_A.
\end{cases} \tag{5a}\]

Thus the exact implementation set, including ties, is
\(\{\tau\geq0:J_\tau=W^*\}\). For a unique A-first optimum with
\(w_B>0\), it is \([0,\min\{b_A,t_B\})\); with \(w_B=0\), it is
\([0,b_A)\). For a unique B-first optimum with \(b_A\leq0\), it is
\([0,b_B)\); with \(b_A>0,w_A=0\), it is \([b_A,b_B)\); with
\(b_A>0,w_A>0\), it is empty under the early-service tie convention. If
both research gains are nonpositive, every nonnegative fee implements
immediate action. The strict-urgency qualification below distinguishes
the empty-set claim from an impossibility under all favorable
alternative tie-breaking. Equation (5a) follows directly from the two
equilibrium branches: a charge can change a choice, but the charge is
not itself a social resource expense.

If \(\ell_A=0\), then \(h_A=d_A=0\) because \(g_A>0\). Early and late
private service can tie. Under favorable alternative tie-breaking, an
efficient B-first schedule with A retained can then be supported at zero
charge. Our robust impossibility result requires strict local urgency,
not an arbitrary unfavorable tie. A change in priority ownership or a
dispatcher that already assigns the efficient order can likewise
implement at zero charge. We are evaluating charges within the specified
decentralized institution, not proving every constant-price mechanism
impossible.

\subsection{6.2 Common deadline}\label{common-deadline}

A common binding deadline leaves \(D=0,1,\) or \(2\) slots for both
cases; it changes neither priors nor priority rights. At zero charge its
equilibrium payoffs are

\[J_0=G,\]

\[J_1=\begin{cases}G+b_A,&b_A>0,\\G+b_B^+,&b_A\leq0,\end{cases}\]

\[J_2=\begin{cases}G+b_A+w_B,&b_A>0,\\G+b_B^+,&b_A\leq0.\end{cases} \tag{5}\]

\textbf{Corollary 1.} If B first is uniquely efficient and \(b_A>0\), no
unpriced common deadline implements the optimum. Conversely, if A first
is efficient, \(D=2\) implements an efficient policy; if no research is
valuable, \(D=0\) does so. Equation (5) gives the exact test including
ties and other cases.

\textbf{Proof.} Under unique B-first optimality,
\(W^*>G+b_A+w_B\geq G+b_A>G\). These three quantities bound the payoffs
of the three deadlines when A has positive research value. The converse
statements follow from (4)-(5). \(\square\)

The deadline can reduce total research while preserving the wrong
service order. This is a different failure from an externally imposed
time price that is simply too low. A deadline with case-specific
priority reassignment is a different intervention and can work.

\subsection{6.3 State-contingent opportunity-loss
charge}\label{state-contingent-opportunity-loss-charge}

Replace the common charge by a posted priority-right rule. If A claims
the first slot, it pays \(u_B=b_B^+-w_B\) once. If it releases the slot,
it pays zero. Deferred experiments and B's use are uncharged. B can
decline research and act; A can retain its second-slot option. The
charge is public, computed before choices from the competing case's
primitives, and does not count as a social expense.

\textbf{Theorem 3 (public-state implementation).} The posted
priority-right rule implements an optimal joint stopping and scheduling
policy for every parameter configuration in the model, with
payoff-equivalent tie choices allowed. If both cases have positive
research value, A chooses first service exactly when \(u_A>u_B\),
otherwise it releases the slot. Waiting and action choices remain
efficient.

\textbf{Proof.} A's payoff gain from claiming is \(b_A-u_B\). Its payoff
gain from releasing to a willing B and choosing its best waiting-or-act
option is \(w_A\). For positive \(b_A,b_B\), the difference is
\(u_A-u_B\). B waits after an A claim if \(w_B>0\) and researches
immediately after a release if \(b_B>0\). Thus the selected route is
exactly the larger route in (4).

If \(b_A\leq0\), its claim cannot improve on releasing and acting; B
then researches if worthwhile. If \(b_B\leq0\), \(u_B=w_B=0\). A claims
if immediate research dominates its waiting option, and otherwise may
take its identical-payoff research option in the final slot. If neither
has positive \(b\), both act. At equality either order has the same
total value. No future charge distorts the second-slot option.
\(\square\)

\textbf{Corollary 2 (when rival-state variation is necessary).} Fix A
with positive b\_A and w\_A. Suppose two public rival states both have
b\_B\textgreater0 and satisfy \(u_B^{L}<u_A<u_B^{H}\). No single
first-slot fee k, independent of the rival state and with uncharged
deferred service, implements the efficient order in both states under
the same tie rule. The opportunity-loss charge does.

\textbf{Proof.} For a constant first-slot fee, A compares \(b_A-k\) with
w\_A, so it claims exactly when \(k<u_A\), independently of B's state.
Efficiency requires claiming in the low state and releasing in the high
state, contradictory strict requirements on k. Posting \(u_B\)
implements each state by Theorem 3. \(\square\)

The rule charges the rival's \emph{loss from deferral}, not its entire
research value. Charging \(b_B^+\) would overstate the externality
whenever \(w_B>0\). Repeating the priority toll on a later slot would
charge for the same lost option twice and could recreate underresearch.
This is precisely why the distinction between a first-slot right and
every unit of research is substantive.

Implementation requires observable case states, enforceable priority
rights and marginal budget incentives. A recommendation that does not
affect behavior is not an implementation. The rule is feasible as an
internal booked priority charge or as equivalent dispatcher assignment
with recorded opportunity cost. It does not require an outside payment.
It is not claimed budget balanced among owners, a truthful auction, or
valid under private misreporting. Dynamic pivot and congestion pricing
precedents remain the appropriate conceptual comparison.

\section{7. Recoverability separates research value from
priority}\label{recoverability-separates-research-value-from-priority}

Consider a case in the interior region \(b_i>0,w_i>0\). Its priority
loss simplifies to

\[u_i=h_ib_i+\ell_i=h_is_iM_i-h_ic_i+s_id_i. \tag{6}\]

\textbf{Theorem 4 (recoverability-priority reversal).} Holding the
prior, signal law, cost, and expiration fixed, within this region

\[\frac{\partial b_i}{\partial r_i}=-S_i[(m_i-\bar m_i)+h_i\bar m_i]\leq0,\]

\[\frac{\partial u_i}{\partial r_i}=h_is_iS_i\bar m_i\geq0. \tag{7}\]

The inequalities are strict when their respective coefficients are
positive. Therefore increasing recoverability can lower a case's
standalone research gain while raising its scheduling priority. Against
a rival with fixed \(u_j\), a priority reversal occurs whenever two
recoverability levels inside this region satisfy
\(u_i(r_0)<u_j<u_i(r_1)\).

\textbf{Proof.} Differentiate \(g_i\) and \(M_i\) in (1): their
derivatives are \(S_im_i\) and \(S_i\bar m_i\). The derivative of
\(b_i\) is \(s_iS_i\bar m_i-S_im_i\). With \(w_i=s_ib_i-\ell_i>0\), (6)
follows algebraically and its derivative yields (7). Theorem 1 converts
crossing losses into a reversed order. \(\square\)

The economic reason is conditional. Recoverability raises the value
preserved by acting in the wrong state. Within the region where research
still occurs, it also raises the expected reward whose availability is
threatened by expiration during an extra waiting slot. The value of
improving the choice falls, but the cost of delaying that
research-and-action option can rise. Once deferred research becomes
unattractive, \(u_i=b_i^+\) and the rising-priority conclusion need not
hold. At \(h_i=0\), the positive derivative in (7) vanishes.

This reduced-form result is not a theorem that reversible trials should
be prioritized. A trial that generates new information after action or
changes the opportunity's hazard requires a different model. Nor is
increased recoverability equivalent to a lower loss from expiration:
expiration destroys the entire still-unresolved opportunity here. If
salvage survives expiration, (1) and the comparative statics change. The
theorem makes those boundaries explicit.

\section{8. Correlated states: the smallest nonseparable
extension}\label{correlated-states-the-smallest-nonseparable-extension}

\subsection{8.1 Information and the contingent joint
policy}\label{information-and-the-contingent-joint-policy}

Keep the same two cases, two slots, one experiment per case,
holding/expiry technology and immediate terminal-action option. Replace
the product prior by a public joint binary distribution
\(\mu=(\mu_{00},\mu_{01},\mu_{10},\mu_{11})\). Conditional on the latent
state vector, signal errors remain independent. Survival is independent
of latent states and signals. Terminal action locks a payoff, but the
realized action outcome is not revealed during the two-slot window;
otherwise correlated cases could learn for free from action outcomes.
Only research signals provide new state information here.

Let \(L\) be the live unresolved cases and \(E\) the cases with an
unused experiment. Observing signal y from research on k updates the
\emph{whole} joint distribution to \(\mu^{k,y}\). Write \(g_i(\mu)\) for
immediate-action value and

\[M_{i|k}(\mu)=E_y g_i(\mu^{k,y}),\qquad
b_k(\mu)=s_kM_{k|k}(\mu)-c_k-d_k-g_k(\mu).\]

The cross-information net gain for a passive case i retained while k is
researched is

\[a_{i|k}(\mu)=s_is_k[M_{i|k}(\mu)-g_i(\mu)]-\ell_i(\mu),
\quad\ell_i(\mu)=d_i+h_ig_i(\mu). \tag{8}\]

The factor \(s_is_k\) matters: i must remain actionable and k must
survive to deliver the signal. If k expires, i still acts on the
unchanged posterior, provided i survives. A previously researched case
can now have \(a_{i|k}>0\) and should not be automatically retired.

\textbf{Theorem 5 (complete correlated two-slot solution).} With one
slot left,

\[V_1(\mu,L,E)=\sum_{i\in L}g_i(\mu)+
\max\left\{0,\max_{k\in L\cap E}
\left[b_k(\mu)+\sum_{i\in L\setminus\{k\}}a_{i|k}(\mu)^+\right]\right\}. \tag{9}\]

If k is selected, retain each other live case i exactly when
\(a_{i|k}>0\); otherwise act on it before the signal. After survival and
signal revelation, choose posterior-optimal terminal actions.

At the initial boundary choose a retained subset \(R\subseteq\{A,B\}\)
and research \(k\in R\) or \(k=\varnothing\). Define
\(E_k=\{A,B\}\setminus\{k\}\) when k is a case, and
\(E_\varnothing=\{A,B\}\). For each survivor subset \(S\subseteq R\) let
\(P_R(S)=\prod_{i\in S}s_i\prod_{i\in R\setminus S}h_i\). Let
\(\mathcal T_k V_1\) be the signal expectation of
\(V_1(\mu^{k,y},S,E_k)\) if \(k\in S\), or \(V_1(\mu,S,E_k)\) otherwise.
Then

\[Q_0(R,k)=\sum_{i\notin R}g_i(\mu)-\sum_{i\in R}d_i
-\mathbf1\{k\ne\varnothing\}c_k
+\sum_{S\subseteq R}P_R(S)\mathcal T_kV_1, \tag{10}\]

\[V_2(\mu)=\max_{R,k}Q_0(R,k). \tag{11}\]

There are exactly eight root choices: stop both; keep only A with
research or idling; keep only B with research or idling; keep both with
A research, B research, or idling. Equations (8)-(11), involving only
finite sums, posterior updates and maxima, solve every initial and
signal-contingent decision; no general queue equation is left unsolved.

\textbf{Proof.} In the final slot either all cases act or one unused
experiment is selected. Its own payoff relative to acting is b\_k. For
another case i, retaining it produces
\(-d_i+s_i[(1-s_k)g_i+s_kM_{i|k}]\), so its incremental gain is (8).
Choose retention separately by its sign; terminal payoffs are additive
and there are at most two cases. Final idling gives no information and
is dominated by immediate action because costs are nonnegative and
rewards nonnegative. This proves (9). For the first slot, outside-R
actions are irrevocable, every retained case pays its holding cost once,
and the selected experiment costs c\_k once. Condition on independent
survival and the signal, if delivered; the remaining optimum is (9). The
eight choices exhaust the initial action set, including idle policies.
This proves (10)-(11) and the contingent policy. \(\square\)

For a product prior, \(M_{i|k}=g_i\) when i differs from k, so every
passive cross gain is nonpositive. Previously researched cases then
retire immediately, and (9)-(11) reduce to Theorem 1. This proves an
exact separable boundary rather than treating correlation as another
scalar delay cost.

\subsection{8.2 Local behavior and a signed incentive
boundary}\label{local-behavior-and-a-signed-incentive-boundary}

The local-priority game extends the earlier one without preventing
cross-learning. In the last slot, if A still has an experiment, it
compares its own charged research gain \(b_A-\tau\) with the gain
\(a_{A|B}^+\) from releasing to B, provided B has an experiment and
\(b_B>\tau\). If A claims, B retains only when \(a_{B|A}>0\). If A
releases, B researches only when its own \(b_B>\tau\), and A retains
only when \(a_{A|B}>0\). If A has no experiment, B has the research
choice. Each owner may act immediately rather than wait. Ties between
research and action favor action; otherwise use the earlier priority
convention. This explicitly solves the last-slot local game.

For completeness the root game is also finite and solved. Let
\(Q_i^L(R,k)\) be case i's component of (10) when its final continuation
is the just-described local game rather than (9), and subtract any own
research charges in computing its private component \(P_i^L\). A
initially chooses claim, release-and-retain, or act-and-release. B's
feasible root choices in those three branches are respectively

\[\mathcal C=\{(\{A\},A),(\{A,B\},A)\},\]
\[\mathcal R=\{(\{A\},\varnothing),(\{A,B\},B),(\{A,B\},\varnothing)\},\]
\[\mathcal X=\{(\varnothing,\varnothing),(\{B\},B),(\{B\},\varnothing)\}.\]

In each branch B selects a maximizer of \(P_B^L\); A then selects the
resulting branch maximizing \(P_A^L\). This is an explicit two-level
maximum over 2, 3, and 3 choices with known finite continuation values.
It specifies off-path actions as well as the equilibrium. Appendix D
proves the backward-induction construction and the code reports all
route payoffs. It reduces to the earlier local equilibrium for
independent states.

\textbf{Theorem 6 (information-credit boundary).} At a last-slot history
with live cases i,j and only i's experiment remaining, joint research is
strictly optimal if and only if \(b_i+a_{j|i}^+>0\). Local research at a
nonnegative common fee requires \(b_i>\tau\). Therefore, if

\[b_i\leq0<b_i+a_{j|i}^+, \tag{12}\]

no nonnegative research charge implements the joint-optimal
continuation. A public credit \(\kappa_i=a_{j|i}^+\) conditional on
conducting the experiment, with no further charge, implements the joint
research-or-action decision. A credit strictly exceeding \(-b_i\) is
necessary for strict local research in the negative-own-gain case; its
infimum is \(-b_i\). These statements concern the last-slot decision,
not a claim that this credit alone implements all root incentives.

\textbf{Proof.} Equation (9) compares exactly those two choices when
only one experiment remains. The passive owner retains iff its own a is
positive. The active owner's charged gain is \(b_i-\tau\), which is
nonpositive under (12). With the credit its gain becomes
\(b_i+a_{j|i}^+\), identical to the joint comparison. Equality can
select action, which is then jointly optimal. Strict local research
requires \(b_i+\kappa_i>0\). \(\square\)

The independent common-fee classification and positive first-slot toll
remain correct on the product-prior subspace; they do not extend by
substituting marginal beliefs into b,w,u. Correlation can create both an
incentive to free-ride on another case's signal and an incentive to stop
socially valuable cross-informing research. A blanket assertion that
local objectives always overacquire evidence would now be false.

A sufficient \emph{full} public-state implementation is available if
unrestricted signed internal transfers and observable eventual resource
payoffs are allowed: give owner i its own resource payoff plus the other
case's resource payoff minus a fixed initial constant. Each owner then
maximizes total firm payoff minus a constant at every decision. Backward
induction through the finite game implements (9)-(11). This elementary
team-payoff alignment is not a new pivot theorem, not budget balanced
between owners, and not a private-information mechanism. It requires
stronger incentive and accounting authority than the one-time toll. The
last-slot credit alone should not be described as implementing the full
root policy.

\subsection{8.3 Exact positive-expiry example and contingent
choices}\label{exact-positive-expiry-example-and-contingent-choices}

Take \(\mu=(.58,.02,.02,.38)\), so both marginal state-1 probabilities
are .4 and all joint states have positive probability. Let
\(S_A=S_B=100\), \(r_A=r_B=0\), \(d_A=d_B=0\), \(h_A=h_B=.01\),
\(q_A=.8,c_A=1\), and \(q_B=.95,c_B=30\). These are arbitrary synthetic
model units.

The unique best root choice retains both and researches A first,
yielding

\[V_2=\frac{249523577}{1562500}=159.69508928.\]

The complete root comparison is:

\begin{longtable}[]{@{}llr@{}}
\toprule\noalign{}
Retained cases & First research & Optimal joint continuation payoff \\
\midrule\noalign{}
\endhead
\bottomrule\noalign{}
\endlastfoot
None & None & 120.000000 \\
A & None & 137.418000 \\
A & A & 138.200000 \\
B & None & 123.409500 \\
B & B & 124.050000 \\
A and B & None & 153.176525 \\
A and B & A & 159.695089 \\
A and B & B & 154.405320 \\
\end{longtable}

If both survive and A's signal is 1, the joint posterior is
\((29/110,1/110,2/55,38/55)\). The posterior marginals are \(p_A=8/11\)
and \(p_B=.7\). B's own last-slot research gain is \(-5.95\), but A's
net gain from being retained to learn B's signal is
\(9401/550=17.09272727\). Thus the joint incremental gain is
\(12257/1100=11.14272727\). Research B and retain A, even though A has
already used its experiment. This is a supported branch with probability
.431244 from the root.

If both survive and A's signal is 0, act on both immediately: the second
experiment is not jointly worthwhile. If A expires but B survives, B's
posterior is unchanged and B researches, with own gain 4.05. If only A
survives, it acts on its observed posterior. If neither survives, there
is no action reward. These six branches are the complete contingent
optimum and are reported with exact probabilities in the reproduction
output.

In the solved local game A instead releases to B initially to obtain
more valuable free cross-information; the total payoff is 154.40532. The
R1 priority toll calculated only from initial marginals is .6405 and
gives the same suboptimal outcome. Independently of this root
comparison, every nonnegative uniform research fee fails to implement
the efficient full policy, since the positive-probability second-slot
branch requires B to research with negative own gain. The unpriced
common-deadline payoffs are 120 for zero slots, 154.22514 for one slot,
and 154.40532 for two slots, all below the joint optimum. Unlike the
independent proof, these comparisons use the correlated local game.

For an exact independence check, replacing the joint prior by
\((.36,.24,.24,.16)\) preserves both marginals and every other
primitive. The resulting joint and local values both equal 141.6095 and
the passive cross-information gain is nonpositive. For a simple boundary
check with perfect positive correlation and no expiry,
\(\mu=(.6,0,0,.4)\) gives a joint value 165.4 versus local 160. These
are three exact small instances, not a simulation campaign or empirical
estimates.

The rational solver and independent physical-outcome enumeration agree
on all eight root values in each instance. The enumerator includes
contingent retention of the already researched case, actual terminal
action success/error, and all positive-probability signal and expiry
paths. There are 63 policy maps in each full-support instance and 30 in
the zero-expiry boundary, 156 total. Final idle policies are
analytically dominated; initial idle policies are enumerated. The new
script reports 37 checks. A separately evaluated full physical tree, not
merely the same Bellman recursion twice, supplies the comparison.

\subsection{8.4 Two public-history credits: complete
implementation}\label{two-public-history-credits-complete-implementation}

The final-slot information credit in Theorem 6 does not resolve A's
initial incentive to free-ride on B's information. We now close that gap
within the \textbf{same} model using only two nonnegative internal
credits. Credit C is paid to A when its initial-slot experiment is
conducted, even if A then expires. Credit \(\kappa\) is paid to B when
its final-slot experiment is conducted after initial research on A, both
cases surviving that first slot, and public signal \(Y_A=1\). Payment
precedes second-slot survival. All other payments are zero, including
every research fee. The announced credits are committed and enter the
owners' quasilinear objectives; they depend on public research/history,
never eventual action rewards. They do not add to firm resource losses.

Keep the synthetic full-support witness's primitives: \(p_A=p_B=.4\),
\(q_A=.8\), \(q_B=.95\), \(c_A=1\), \(c_B=30\), \(h_A=h_B=.01\),
\(S_A=S_B=100\), and \(d_i=r_i=0\). Vary only

\[\mu(\delta)=(.36+\delta,.24-\delta,.24-\delta,.16+\delta),
\qquad -.16\leq\delta\leq .24.\]

We prove the following result on the exact, nonempty closed interval

\[I=[17/100,23/100]. \tag{13}\]

This is a certified interval around .22, not an asserted maximal region
or a classification of the entire admissible prior range. Independence
remains at \(\delta=0\), outside I, and retains Theorems 1-4's
appropriate scope.

The efficient contingent policy \(\pi^*\) on I retains both cases and
researches A initially. If both survive and \(Y_A=1\), it retains A and
researches B; if \(Y_A=0\), both act immediately. If A expires and B
survives, B researches without a credit. If only A survives, A acts on
its own signal. Neither survivor means no terminal reward. After the
second experiment, survivors act using exactly their available public
signals. With \(s=99/100\), its resource components are

\[X_A(\delta)=\frac{43226820450\delta+33275667541}{500000000},\]
\[X_B(\delta)=\frac{297(495000\delta+515039)}{2500000}. \tag{14}\]

Define the credit-event probability, two switch points, and B's
continuation deficit by

\[p=\frac{107811}{250000}=s^2\Pr(Y_A=1),\quad
\delta_0=\frac{7831}{45000},\quad
\delta_1=\frac{5291}{30000},\quad
K(\delta)=\frac{30000\delta-5291}{220}. \tag{15}\]

Here p is \textbf{not} multiplied by second-slot survival: B is paid on
conducting research. Define A's required initial compensation

\[F(\delta)=\begin{cases}
(877679550\delta+171175559)/500000000,&\delta\leq\delta_0,\\
(44982179550\delta-7503987541)/500000000,&\delta\geq\delta_0.
\end{cases} \tag{16}\]

The pieces agree at \(\delta_0\) and F is strictly positive throughout
I.

\textbf{Theorem 7 (two-credit implementation and initial toll failure).}
On I the root action and supported continuation of \(\pi^*\) are
strictly efficient. Under the stipulated zero-gain action tie and
positive initial-priority tie conventions, the two-credit local game
implements its entire contingent policy if and only if

\[C\geq F(\delta),\qquad \kappa\geq0,\qquad
\kappa>K(\delta). \tag{17}\]

Strict implementation at the initial comparison and every supported
discretionary research/retention decision additionally requires
\(C>F(\delta)\). A positive tie between initial claim and release
selects A's initial claim; at \(\kappa=K\geq0\), B instead selects
action and implementation fails. The strict-implementation statement
concerns the efficient outcome and its supported decisions, not
elimination of all possible indifferences in off-path subgames.

Replacing C by an initial A-claim \textbf{toll} \(T\geq0\), while
retaining the same B-only continuation credit and zero other payments,
cannot implement \(\pi^*\) anywhere on I. This is a claim about that
restricted instrument class, not arbitrary taxes, subsidies, or
budget-balanced mechanisms.

\textbf{Proof.} The joint-policy certificate and all off-path
continuation calculations are given in Appendix E. Conditional on both
cases alive after \(Y_A=1\), B's own research gain is \(-K\), and A's
net passive cross gain is

\[a_{A|B}=\frac{4410450\delta-594259}{22000}>0.\]

Their joint increment is \((1410450\delta-65159)/22000>0\) on I. B
conducts research exactly when \(\kappa>K\). A then strictly retains.
After \(Y_A=0\), B's uncredited research gain is
\(-3(5000\delta-189)/140<0\); neither research nor idle continuation is
chosen. On the B-only survival branch, its uncredited gain is
\(81/20>0\). All these decisions use the actual local continuation, not
a prescribed dispatcher.

When \(\kappa>K\), B's root private payoff from retention after an A
claim is \(X_B+p\kappa>60\), so it strictly retains instead of acting.
If A releases and remains unresolved, B strictly researches initially:
its own value is \(1281/20=64.05\), versus \(126819/2000=63.4095\) from
waiting and 60 from acting. These comparisons include the off-path
no-signal final game. The same choices hold if A acts and releases. A's
release-and-retain resource payoff is

\[R_A(\delta)=51.54336+176.418\delta+
\frac{401841}{1000000}\left[\frac{7831-45000\delta}{205}\right]^+ . \tag{18}\]

This includes A's own later research when beneficial after B's first
signal, and research if B expires without delivering a signal. Direct
subtraction gives \(R_A-X_A=F>0\). A's three resulting private payoffs
are \(X_A+C\), \(R_A\), and 60. Since \(R_A>60\), the initial comparison
is exactly \(C\geq F\), with strictness as stated. If \(\kappa\leq K\),
a positive-probability efficient continuation is absent regardless of C;
necessity follows. Conversely all the displayed strict continuation
comparisons and the root comparison imply the complete target policy,
proving sufficiency. With a toll the claim payoff is \(X_A-T<R_A\)
whenever B's target continuation is induced; if not induced the
continuation already fails. This proves the restricted toll
impossibility. \(\square\)

The threshold at \(\delta_0\simeq .17402222\) comes from A's uncredited
own experiment after B's initial signal 1, an off-path continuation that
changes its release payoff. The threshold at
\(\delta_1\simeq .17636667\) comes from B's credited on-path experiment
changing from privately profitable to nonpositive. Both are inside I.
The full game therefore yields piecewise initial and continuation
requirements even though its efficient policy is unchanged.

\subsection{8.5 Funding frontiers, attainment, and exact
witness}\label{funding-frontiers-attainment-and-exact-witness}

Let \(k(\delta)=\max\{0,K(\delta)\}\). A gross expected-credit
constraint \(\bar E\) limits expected transfers, while gross pathwise
capacity \(\bar H\) must fund the largest possible payment total. These
are institutional limits on credits, not welfare-cost terms. Under
implementation, A is always paid C and B is paid \(\kappa\) on an event
of probability p.~Therefore

\[E=C+p\kappa,\qquad H=C+\kappa. \tag{19}\]

\textbf{Theorem 8 (restricted funding frontier).} For weak
implementation selected by the stated ties, the two infima are

\[E_{\inf}(\delta)=F(\delta)+p k(\delta),\qquad
H_{\inf}(\delta)=F(\delta)+k(\delta). \tag{20}\]

If \(\delta<\delta_1\), both are attained together at
\((C,\kappa)=(F,0)\). If \(\delta\geq\delta_1\), neither is attained.
Strict supported implementation has the same two infima and attains
neither at any \(\delta\in I\).

For a weak implementation, the exact budget-feasibility conditions are

\begin{longtable}[]{@{}lll@{}}
\toprule\noalign{}
Prior region & Expected cap only & Pathwise cap only \\
\midrule\noalign{}
\endhead
\bottomrule\noalign{}
\endlastfoot
\(\delta<\delta_1\) & \(\bar E\geq F\) & \(\bar H\geq F\) \\
\(\delta\geq\delta_1\) & \(\bar E>F+pk\) & \(\bar H>F+k\) \\
\end{longtable}

For strict supported implementation, both respective conditions are
strict throughout I. If both budgets are imposed, \textbf{both}
corresponding inequalities are necessary and sufficient. Expected and
pathwise caps cannot be substituted for each other. These statements
completely classify funding inside the two-credit class, without
optimizing over other contracts.

\textbf{Proof.} Equation (17) is the complete feasible payment set. Both
objectives in (19) increase strictly in each payment since \(0<p<1\).
When K is negative, \((F,0)\) is feasible and minimizes both. At K zero
or positive, B's strict research incentive requires \(\kappa>k\);
approach but do not reach k while C equals F. Strict root implementation
also requires C above F. These facts prove the infima and attainment
statements. For simultaneous caps with positive slack, choose
sufficiently small nonnegative increments above F and k; the expected
increment is \(\epsilon_C+p\epsilon_\kappa\), and the pathwise increment
is \(\epsilon_C+\epsilon_\kappa\). They fit both caps. If K is negative
and only weak implementation is required, the attained pair fits both
weak caps directly. Necessity follows from monotonicity. \(\square\)

At the preserved witness \(\delta=11/50\), the exact values are

\[X_A=\frac{1069639201}{12500000}=85.57113608,\quad
R_A=\frac{2258883}{25000}=90.35532,\]
\[F=\frac{59802299}{12500000}=4.78418392,\qquad K=\frac{119}{20}=5.95,\]
\[E_{\inf}=\frac{183752143}{25000000}=7.35008572,\quad
H_{\inf}=\frac{134177299}{12500000}=10.73418392. \tag{21}\]

Neither weak nor strict implementation attains these witness infima: B
chooses action at \(\kappa=5.95\). At that limiting value B's root
retained-route private value is \(X_B+pK=76.689855>60\), verifying that
retention is slack. For example \(C=F+.01\) and \(\kappa=K+.01\)
implement strictly, with \(E=7.36439816\) and \(H=10.75418392\). The
firm's physical payoff remains 159.69508928; adding or subtracting the
credits from it would conflate internal funding with resource
accounting.

The exact endpoint certificate and witness computations enumerate
physical latent states, signals and survival with rational arithmetic,
independently of the joint-value solver. They verify 63 contingent
schedule/retention maps at each of two interval endpoints and the
retained witness. This is a proof certificate and accounting check, not
a parameter grid or new simulation campaign. The 101 checks include all
root alternatives, complete credited local-game choices, both binding
incentive endpoints, piece continuity, and the original fractions.

\section{9. Exact synthetic
verification}\label{exact-synthetic-verification}

\subsection{9.1 Independent computation and
counting}\label{independent-computation-and-counting}

All reported numbers are synthetic model units. The retained
independent-case script \texttt{joint\_analysis.py} computes (1)-(5) and
the posted policy. Independently, a second evaluator enumerates every
root choice of retained subset and active case or idling, every feasible
signal-contingent final-slot research map, both latent states, both
signal outcomes, and all positive-probability survival paths. It values
realized correct and incorrect terminal actions and actual costs
directly; it does not use the \(b,w,u\) formulas in the evaluator.

With strictly interior hazards the full two-case tree contains 35
distinct scheduling policies in this representation: immediate action on
both; one-case research; one-case idling with terminal choice; two-case
idling with survivor-contingent scheduling; and research on either case
with signal-contingent final allocation. When zero-probability branches
are removed the count can be smaller. Terminal actions after each
available signal use Bayes-optimal classification; all possible research
schedules are enumerated. The separate posterior computation does not
import the closed-form scheduling optimum.

For the six named scenarios the output contains 193 policy rows,
including 35 for each interior-hazard example and 18 when both hazards
vanish. The script also checks additional parameter configurations and
the policy boundaries; 250 independent-case checks pass. These checks
validate implementation in this finite model, not external validity or
originality. The earlier single-case script separately passes 12,002
checks and remains in the reproduction package.

\subsection{9.2 Congestion, charge failure, and deadline
failure}\label{congestion-charge-failure-and-deadline-failure}

Let both cases have \(p=.5,S=100,r=0,c=1,d=0\). Let A have
\(q_A=.7,h_A=.02\), and B have \(q_B=.9,h_B=.20\). Then

\[b_A=17.6,\quad w_A=16.248,\quad u_A=1.352,\]
\[b_B=21,\quad w_B=6.8,\quad u_B=14.2.\]

The joint optimum researches B first, retains A, and obtains
\(100+21+16.248=137.248\). Local priority gives A first and B later,
yielding \(124.4\). A one-slot common deadline yields \(117.6\); an
immediate deadline yields \(100\). For every common charge
\(\tau<17.6\), A still takes the first slot; for \(\tau\geq17.6\), A
does not perform the valuable deferred experiment. Thus no common charge
implements the optimum, a conclusion from the theorem rather than a
finite tax grid. The posted priority charge is 14.2. A compares
\(17.6-14.2=3.4\) with its retained value 16.248 and releases the slot,
implementing the optimum.

Changing only B's prior to .705 in this example gives \(g_B=70.5\),
\(b_B=.5\), \(w_B=0\), and \(u_B=.5\). A should then be first, with
total payoff 138.1. A's own primitives are unchanged. The pair of B
priors .5 and .705 straddles A's priority loss, so even a constant
first-slot fee cannot implement both states. This demonstrates why the
correction must respond to the competing belief rather than merely to
the focal analyst's own research cost.

\subsection{9.3 Priority reversal and positive implementation
examples}\label{priority-reversal-and-positive-implementation-examples}

For the reversal, keep A as above and set B to \(q_B=.7,h_B=.022\), with
the other primitives unchanged. Initially \(u_A=1.352<u_B=1.48412\), so
B is first. Increase A's recoverable share to .5. Its research gain
falls from 17.6 to 7.3, but its priority loss rises to 1.646; both
research-retention inequalities remain strict. A is now first. The
payoff levels differ because recoverability changes the primitive
rewards; the comparison identifies an ordering reversal, not a
treatment-effect estimate.

If both hazards and holding costs are zero, \(w_i=b_i^+\) and all
deferral losses vanish. Both positive-value experiments can be performed
in either order at no loss, and a zero common charge works. If A's
acquisition cost rises to 30 in the congestion example, \(b_A=-11.4\)
and it acts immediately; B researches first at zero common charge, also
efficiently. These examples prevent an unconditional claim that constant
charges or existing priority rights always fail.

\section{10. Welfare measurement and prospective queue
experiment}\label{welfare-measurement-and-prospective-queue-experiment}

The proposed empirical unit changes with the shared-resource mechanism.
The primary unit is a pre-registered analyst-queue episode containing a
frozen cohort of eligible decisions and a specified capacity window, not
a completed individual case. The estimand is the intention-to-treat
effect of an offered opportunity-loss scheduling aid on total payoff
over all cases in that episode. Per-case effects and cross-case
spillovers are secondary. No firm data, partner agreement, recruitment,
treatment assignments, or implemented experiment exist.

The protocol crosses an incremental evidence packet with a queue-level
priority aid, keeping analyst capacity and the economic alternatives
comparable. The aid explicitly records immediate action, deferred
research, expiration, and the foregone value of the competing case. It
compares a local-order baseline with opportunity-cost scheduling. It
does not purport to identify the behavior of real budget charges if only
a recommendation is offered. Any charge experiment would require actual
internal authority and an additional registered arm. A common-deadline
arm is neither necessary nor currently authorized. If cases have
correlated underlying conditions, record cross-case prior links,
posterior revisions after each public signal, and decisions to retain an
already researched case. The independent u-score is not applied to such
cases as if a marginal-prior update were sufficient; the two-slot
correlated formula or an explicitly limited recommendation must be used.

Randomization must occur at the entire shared analyst pool or
independent pool-period level, with buffer periods and all displaced
work tracked. Randomizing individual cases within the same queue would
expose controls to treatment-induced delay and cannot identify an
isolated case effect without an interference design. The episode cohort
includes rejected, expired, and unresolved cases. Primary payoffs
include every case and every worker, not just those served before the
measurement window ends.

Double counting is a central risk. Actual action margin, research cost,
residual holding cost, expiration consequence, and correction/rework
expense must have exclusive ledger categories. The theoretical \(u_j\)
already summarizes the rival's modeled forgone option and is not added
to its observed payoff loss. Internal tolls are transfers. Uncertain
costs produce conditional valuation bounds rather than arbitrary
monetary estimates. Confidence intervals must account for both
assignment and partial identification.

If only \(d_i\) is interval valued, every comparison remains conditional
on priors, signal laws, terminal rewards, expiration, and feasible
actions. If those are also uncertain, the entire set of candidate joint
policy values must be recomputed over a documented joint uncertainty
set. A policy is certified only if its worst comparison is nonnegative
throughout that set. No unconditional model-robust claim follows from
bounding delay prices alone.

The detailed prospective protocol contains queue and case dictionaries,
inclusion rules, interference assumptions, missing-outcome bounds, and
illustrative precision calculations. Public BPI Challenge 2019
documentation provides procurement-process context; its repository file
and license were not verified and no rows were analyzed. It cannot
identify the proposed mechanism without resource rosters, evidence
requests, complete intake, and linked economic outcomes.

\section{11. Conclusion}\label{limitations-and-conclusion}

The scarce resource is research, not terminal authorization; action
consumes no shared capacity. There are exactly two cases, one experiment
per case, and two slots. Sections 4-7 assume independent states; Section
8 solves the explicitly correlated counterpart. Neither model includes
more experiments per case, arrivals, heterogeneous durations, endogenous
effort, private reporting, or commitment problems. Terminal outcomes do
not reveal latent states during the two-slot decision window; only
public research signals provide new state information. We do not present
an unsolved general queue as a completed result.

The bounded independent-case contribution is the recoverability/priority
derivative and the exact common-charge implementation set, not the
scheduling interchange or public externality charge themselves. Under
correlation, the first signal changes the second case's research and
retention values. A supported branch can require research despite its
negative own-case gain, so the positive-toll-only extension fails. The
full contingent joint policy and finite local game are solved. Theorems
7-8 close the initial incentive gap for the two specified credits, with
exact piecewise thresholds and expected/pathwise funding limits. Their
nonnegative-payment constraint and public-history support are essential
scope conditions.

These results are stronger than a separable benefit/time ratio but
remain applications of established scheduling and externality
principles. Bergemann-Välimäki provides a much broader mechanism-design
framework, and expiring/restless scheduling is a substantial existing
literature. The strongest remaining objection is that the small
public-state implementation characterization and comparative static may
still be too modest for a leading general-interest journal. No amount of
synthetic verification establishes empirical relevance or priority. The
two-credit frontier is a complete restricted implementation result, but
uses standard incentive comparisons and Bayesian backward induction and
does not certify a substantial originality claim. It does not solve
private reports, optimal choice of payment support, budget balance,
endogenous funding, or arbitrary correlations outside the certified
interval. No empirical evidence resolves practical importance. This is
an integrated revised theoretical paper with a separate prospective
study, not a claim of scientific acceptance.


{
\small
\setlength{\parindent}{0pt}
\section*{References}\label{references}

Adiri, Igal, and Uri Yechiali. 1974. ``Optimal Priority-Purchasing and
Pricing Decisions in Nonmonopoly and Monopoly Queues.'' \emph{Operations
Research} 22(5): 1051-1066. \url{https://doi.org/10.1287/opre.22.5.1051}.

Bergemann, Dirk, and Juuso Välimäki. 2010. ``The Dynamic Pivot
Mechanism.'' \emph{Econometrica} 78(2): 771-789.
\url{https://doi.org/10.3982/ECTA7260}.

Campbell, Arthur, Florian Ederer, and Johannes Spinnewijn. 2014. ``Delay
and Deadlines: Freeriding and Information Revelation in Partnerships.''
\emph{American Economic Journal: Microeconomics} 6(2): 163-204.
\url{https://doi.org/10.1257/mic.6.2.163}.

Escobar, Juan F., and Qiaoxi Zhang. 2021. ``Delegating Learning.''
\emph{Theoretical Economics} 16(2): 571-603.
\url{https://doi.org/10.3982/TE4191}.

Eyring, Henry, Patrick J. Ferguson, and Sebastian Koppers. 2021. ``Less
Information, More Comparison, and Better Performance: Evidence from a
Field Experiment.'' \emph{Journal of Accounting Research} 59(2):
657-711. \url{https://doi.org/10.1111/1475-679X.12362}.

Fudenberg, Drew, Philipp Strack, and Tomasz Strzalecki. 2018. ``Speed,
Accuracy, and the Optimal Timing of Choices.'' \emph{American Economic
Review} 108(12): 3651-3684. \url{https://doi.org/10.1257/aer.20150742}.

Gittins, John C. 1979. ``Bandit Processes and Dynamic Allocation
Indices.'' \emph{Journal of the Royal Statistical Society, Series B}
41(2): 148-164. \url{https://doi.org/10.1111/j.2517-6161.1979.tb01068.x}.

Henry, Emeric, and Marco Ottaviani. 2019. ``Research and the Approval
Process: The Organization of Persuasion.'' \emph{American Economic
Review} 109(3): 911-955. \url{https://doi.org/10.1257/aer.20171919}.

Klein, Nicolas, and Sven Rady. 2011. ``Negatively Correlated Bandits.''
\emph{Review of Economic Studies} 78(2): 693-732.
\url{https://doi.org/10.1093/restud/rdq025}. Read version: author-hosted
manuscript dated June 23, 2010,
\url{https://nicolasklein.com/data/documents/Klein}\_Rady\_CB.pdf; publisher
metadata verifies the later journal version, whose full text was not
read.

Moroni, Sofia. 2022. ``Experimentation in Organizations.''
\emph{Theoretical Economics} 17(3): 1403-1450.
\url{https://doi.org/10.3982/TE3069}. Published introduction and model
inspected.

Righter, Rhonda, and J. George Shanthikumar. 1998. ``Independently
Expiring Multiarmed Bandits.'' \emph{Probability in the Engineering and
Informational Sciences} 12(4): 453-468.
\url{https://doi.org/10.1017/S0269964800005325}.

Weber, Richard R. 1992. ``On the Gittins Index for Multiarmed Bandits.''
\emph{Annals of Applied Probability} 2(4): 1024-1033.
\url{https://doi.org/10.1214/aoap/1177005588}.

Weber, Richard R., and Gideon Weiss. 1990. ``On an Index Policy for
Restless Bandits.'' \emph{Journal of Applied Probability} 27(3):
637-648. \url{https://doi.org/10.2307/3214547}.

Whittle, Peter. 1988. ``Restless Bandits: Activity Allocation in a
Changing World.'' \emph{Journal of Applied Probability} 25(A): 287-298.
\url{https://doi.org/10.2307/3214163}.

Xiao, Yangge, Zhenyu Hu, and Shouqiang Wang. 2026. ``Information Design
of a Delegated Search.'' \emph{Management Science}, Articles in Advance,
published online May 14. \url{https://doi.org/10.1287/mnsc.2024.05925}.

Xu, Renyuan, Thaleia Zariphopoulou, and Luhao Zhang. 2026. ``Decision
Making under Costly Sequential Information Acquisition: the Paradigm of
Reversible and Irreversible Decisions.'' arXiv:2401.00569v4, January 6.
Initial submission December 31, 2023. \url{https://arxiv.org/abs/2401.00569}.
Working-paper version read.

Yu, Zhe, Yunjian Xu, and Lang Tong. 2017. ``Deadline Scheduling as
Restless Bandits.'' arXiv:1610.00399v5, July 6.
\url{https://arxiv.org/abs/1610.00399}. Preprint read. Later publication:
\emph{IEEE Transactions on Automatic Control} 63(8), August 2018, DOI
\url{https://doi.org/10.1109/TAC.2018.2807924}; journal existence and issue
verified from the author's Cornell bibliography, journal full text not
read. No assertion of unchanged theorem numbering between versions.

}
\clearpage
\section{Appendix A. Shared-resource model: complete derivations and
verification}\label{appendix-a.-shared-resource-model-complete-derivations-and-verification}

\subsection{A.1 Binary signal and posterior
actions}\label{a.1-binary-signal-and-posterior-actions}

For case i with prior p and accuracy q, the joint probabilities of state
1 and signal 1, state 0 and signal 1 are pq and (1-p)(1-q). The joint
probabilities for signal 0 are p(1-q) and (1-p)q. Bayes-optimal expected
classification accuracy is

\[a(p,q)=\max\{pq,(1-p)(1-q)\}+\max\{p(1-q),(1-p)q\}.\]

Thus \(\bar m=1-a(p,q)\) and \(M=S[r+(1-r)a(p,q)]\). At any
zero-probability signal branch, the posterior action is irrelevant to
the expected payoff. At a posterior tie either terminal action has the
same value. The enumeration uses a fixed state-1 tie convention. The
formula applies to asymmetric priors without presuming that the signal
always overturns the prior-optimal action.

\subsection{A.2 Deferred value and no double
counting}\label{a.2-deferred-value-and-no-double-counting}

If the case acts now, it receives g. If it waits one slot before
deciding whether to research, it incurs d immediately; expiration
destroys the opportunity with probability h. Conditional on survival the
last-slot optimal value is \(g+b^+\). Waiting's incremental payoff is

\[-d+s(g+b^+)-g=-\ell+sb^+.\]

The best of action now and this option gives \([-\ell+sb^+]^+\). Since
\(\ell\geq0\), this equals \([sb-\ell]^+\) for every real b, which is w.
In particular, w\textgreater0 implies b\textgreater0. The expression
includes the cost of the later experiment through b, the first waiting
slot's cost through ell, and two distinct opportunities for expiration
through s and b. It does not charge a separate hypothetical missed
opportunity in addition to those losses.

For i first and j waiting, a direct physical-outcome calculation gives

\[s_iM_i-c_i-d_i-d_j+s_j(s_jM_j-c_j-d_j).\]

Subtract \(g_i+g_j\) and rearrange to \(b_i+s_jb_j-\ell_j\). If waiting
j is not profitable, immediately act on j and obtain \(G+b_i\). Taking
their maximum yields \(G+b_i+w_j\). This identity is checked by the
independent latent-outcome evaluator, not assumed in that evaluator.

\subsection{A.3 Exhaustiveness of Theorem
1}\label{a.3-exhaustiveness-of-theorem-1}

At the initial boundary there are four retained subsets, including the
empty one. For each retained subset R, the active choice is one case in
R or no research. Cases outside R act immediately. If some case is
active, its experiment uses its only research opportunity; if it
survives, action immediately after its signal weakly dominates further
waiting because no additional own-state information or reward becomes
available. The independent state assumption means its public signal
cannot improve the remaining case's posterior. Therefore the only
remaining decision is whether an unserved case was retained and, if it
survived, whether to research it in the last slot. This gives the two
research-first candidates in Theorem 1, with optional immediate
retirement of the unserved case embedded in w.

If no case is active initially, fix the retained subset R. Its expected
value relative to immediate action on all cases is

\[-\sum_{i\in R}\ell_i+E\max\{0,\{b_i:i\in R\text{ survives}\}\}.\]

The expectation is bounded above by \(\max(0,b_A,b_B)\). Immediate
research on a positive maximizer, if one exists, followed by immediate
action on the other case attains that upper bound without incurring the
idle waiting losses. If none exists, stopping both dominates idling.
This proves exhaustion even for policies that idle while retaining only
one case. Signal-dependent final research choices cannot improve on w,
because the other case's independent signal changes neither b nor the
feasible final choices.

For both b positive, order comparison is exactly

\[[b_A+w_B]-[b_B+w_A]=u_A-u_B.\]

When b is nonpositive, that case's w is zero and it never needs
research; Theorem 1's maximum also captures these corners. A case with
h=0,d=0 can be deferred without loss. This yields payoff ties, not a
failure of the formula.

\subsection{A.4 Sequential equilibrium with a common
charge}\label{a.4-sequential-equilibrium-with-a-common-charge}

The game is complete information. At the final slot each surviving case
researches only if its charged gain b-tau is strictly positive. At the
initial boundary A moves first. If A researches now, B's best wait gain
is \([s_B(b_B-\tau)-\ell_B]^+\) and its alternative is immediate action.
If A releases, B researches immediately whenever \(b_B-\tau>0\);
otherwise it acts. In the latter event A cannot gain by idling instead
of taking its own available experiment immediately. If A releases to a
researching B, its best waiting gain is \([s_A(b_A-\tau)-\ell_A]^+\).

For b\_A-tau\textgreater0, its own immediate gain minus the untruncated
deferred gain is

\[h_A(b_A-\tau)+\ell_A\geq0.\]

Truncating a negative deferred gain at zero preserves weak preference
for immediate research. It is strict when ell\_A\textgreater0. If both
h\_A and ell\_A are zero, early and late gains tie and the stated
early-service convention selects the first slot. For
b\_A-tau\textless=0, immediate action weakly dominates both research and
waiting; the research-action tie convention selects action. Backward
induction therefore yields Theorem 2, including off-path choices.

The actual firm payoff for b\_A\textgreater tau is

\[G+b_A+\mathbf 1\{s_B(b_B-\tau)>\ell_B\}\,w_B.\]

The charge affects the private choice but is not subtracted again from
social payoff. For b\_A\textless=tau it is
\(G+\mathbf 1\{b_B>\tau\}b_B\). These formulas permit evaluation of all
tau, not just sampled values.

\subsection{A.5 Full implementation classification and tie
boundaries}\label{a.5-full-implementation-classification-and-tie-boundaries}

Assume first that the efficient first research assignment is unique. If
A first is optimal, b\_A\textgreater0 and tau=0 reproduces its efficient
retention decision for B. If B first is optimal and b\_A\textless=0,
tau=0 implements it. If neither case has positive b, acting immediately
is efficient and tau=0 implements it.

The remaining case has unique B first and b\_A\textgreater0. If w\_A=0,
strict optimality implies b\_B\textgreater b\_A+w\_B\textgreater=b\_A. A
releases and stops for tau\textgreater=b\_A, while B researches for
tau\textless b\_B; hence {[}b\_A,b\_B) is sufficient and, under the
early-service convention, necessary. If w\_A\textgreater0,
tau\textless b\_A preserves A's first claim and tau\textgreater=b\_A
destroys its later research incentive. Thus no tau works under that
convention. With ell\_A\textgreater0 this conclusion is independent of
first-slot tie-breaking because A strictly prefers early service for
every tau\textless b\_A.

If ell\_A=0, h\_A=d\_A=0 because g\_A\textgreater0. For
tau\textless b\_A A can be indifferent between early and late service.
Selecting release in that tie can support the efficient B-first route.
The stronger claim that \emph{no equilibrium with any tie-breaking}
implements the optimum would therefore be false. We do not make it. At
ties in the joint optimum, implementation requires attaining any
joint-optimal policy, not a predetermined labeled ordering. The formulas
in the main text directly test such cases.

A common deadline can leave zero, one, or two slots because the model is
discrete. Equation (5) is the exhaustive equilibrium payoff list for an
unpriced unchanged priority game. A strict B-first optimum with
b\_A\textgreater0 exceeds all three. If deadline design includes
changing priority rights, case-specific deadlines, or contingent
exceptions, it is outside this impossibility statement.

\subsection{A.6 Opportunity-loss pricing and
transfers}\label{a.6-opportunity-loss-pricing-and-transfers}

Corollary 2 follows because a fixed first-slot fee k induces the same
comparison \(b_A-k\) versus w\_A in every rival state. With A fixed and
both rivals willing to use a released slot, low and high rival deferral
losses that straddle u\_A require opposite strict inequalities on k. In
the synthetic belief example, changing only B's prior from .5 to .705
changes its loss from 14.2 to .5, straddling u\_A=1.352.

The posted first-slot charge for A is \(u_B=b_B^+-w_B\). If both cases
have positive b, A compares its early gain \(b_A-u_B\) with its
release-and-retain gain w\_A. The difference is u\_A-u\_B. B's
first-service and deferred-service choices are uncharged, so b\_B and
w\_B are undistorted. The resulting joint policy selects exactly a
maximizing route in Theorem 1. The same argument applies if A
immediately acts after release because w\_A=0.

If b\_B\textless=0, u\_B=w\_B=0, B immediately acts. A's early versus
deferred comparison is b\_A versus w\_A. If these tie at a positive
value, A has h\_A=d\_A=0, so waiting and using the otherwise idle final
slot is payoff equivalent. If b\_A\textless=0 it cannot obtain positive
gain by claiming or waiting. Both nonpositive cases stop. This proves
all corners of Theorem 3.

The charge equals the change in B's optimal incremental payoff between
being offered first service and being offered deferred service. It is
not B's gross reward or all of its evidence value. A central budget can
retain it. The implementation claim assumes marginal booked transfers
enter objectives and cannot be evaded; it does not establish budget
balance, individual participation relative to every alternative
institutional entitlement, or private-type revelation. The general
dynamic-pivot literature is therefore a precedent, not a theorem proved
anew here.

\subsection{A.7 Comparative statics and a reversed
order}\label{a.7-comparative-statics-and-a-reversed-order}

Let delta=m-mbar\textgreater=0. At fixed signal law,

\[\partial_r b=sS\bar m-Sm=-S(\delta+h\bar m).\]

Within b\textgreater0,w\textgreater0,

\[u=b-(sb-\ell)=hb+\ell=hsM-hc+sd,\]

so \(\partial_r u=hsS\bar m\). At p=.5,q=.7,S=100,c=1,d=0,h=.02,
b(r)=17.6-20.6r and u(r)=1.352+.588r on this interior segment. At r=.5,
b=7.3,w=5.654,u=1.646. A fixed rival u=1.48412 is crossed at r
approximately .22469, inside the region where deferred research remains
profitable. The exact crossing is (1.48412-1.352)/.588. All reported
endpoints and ordering are in the generated output.

For other parameters within this region,
\(\partial_S u=hs[1-(1-r)\bar m]>0\) when h\textgreater0, whereas
\(\partial_S b=s[1-(1-r)\bar m]-[1-(1-r)m]\) may have either sign.
Increasing signal informativeness lowers mbar and raises M; its effect
on b and, in this interior, u is nonnegative. The hazard derivative of u
is \((1-2h)M-c-d\), which need not be positive. Near a stopping
boundary, positive-part kinks require evaluating the piecewise formula,
not extrapolating derivatives. These statements do not extend without
modification to trials, state-dependent expiration, correlated cases, or
multiple remaining signals.

\subsection{A.8 Independent enumeration: policies and physical
paths}\label{a.8-independent-enumeration-policies-and-physical-paths}

The enumerator and closed-form solver share primitive inputs only. It
enumerates theta\_A,theta\_B; one potential binary signal for each case;
and two potential survival coins per case. Unused draws integrate out.
Actual action success or error is evaluated against the latent state;
costs are charged only when a case is retained or researched. A case
stopped at intake earns its actual prior-optimal action payoff on that
latent path, not zero. A case expiring before action receives zero
terminal reward.

For strictly interior hazards, the eight root actions are: stop both;
retain A with active A or no research; retain B with active B or no
research; retain both with A active, B active, or no research. Their
feasible final-slot contingent policies number 1,1,2,1,2,8,8,12,
respectively, totaling 35. For root A-active with both retained, A has
three public outcomes (expiry, signal 0, signal 1); for each, if B
survives the policy can act or research B. Thus there are 2\^{}3=8 maps,
even though independence ultimately makes their optimal choices the same
across A's outcomes. For root idle-both, the both-survive node allows
three final assignments and each one-survive node two; the
neither-survive node has one, giving 3\emph{2}2=12. Zero-probability
nodes are removed, reducing counts in boundary cases.

All positive-probability latent/signal/hazard paths are integrated for
each map. The maximum is compared with Theorem 1. The code also compares
the posted rule with the optimum over a deterministic parameter grid,
verifies local payoff does not exceed it, and checks the charge-failure
and recoverability examples. New validation output reports 250 checks;
these are computational tests, not additional formal results or
estimated observations. The primitive grid is deliberately heterogeneous
and is not a calibration. Complete named-scenario policy rows are
supplied in CSV, including dominated idle policies.

\subsection{A.9 Conditional valuation
bounds}\label{a.9-conditional-valuation-bounds}

For an implemented fixed joint policy z, expected total payoff can be
written \(A_z-\sum_i d_iT_{iz}\), where A includes terminal rewards and
acquisition expenses, and T\_i is expected unresolved case-period
exposure. Expiration is already inside A and T. Holding intervals
\(d_i\in[l_i,u_i]\) give endpoint bounds for each fixed-policy contrast
by selecting endpoints according to the signs of the exposure
differences. Do not add u\_i as another resource expense.

If priors, precision, stakes, recoverability, or hazards are uncertain,
these quantities and even the optimal policy change. A cost-only
endpoint calculation holds them fixed. For a full uncertainty set Omega,
a proposed fixed policy is certified over a comparator set only if the
minimum of each payoff contrast over Omega is nonnegative. Without a
justified Omega and actual parameter restrictions there is no claim of
model robustness. The paper implements synthetic point-primitives only.

\section{Appendix B. Retained separable benchmark and previous
checks}\label{appendix-b.-retained-separable-benchmark-and-previous-checks}

\subsection{B.1 Existence, admissibility, and the cost
frontier}\label{b.1-existence-admissibility-and-the-cost-frontier}

The finite history tree has finitely many actions at each node. A
deterministic strategy assigns stop-and-act or an experiment to each
live nonterminal node. Backward induction gives an optimal deterministic
strategy for every cost price \(b\). Randomized strategies are mixtures
and cannot exceed the best deterministic expected payoff. All rewards
and costs are bounded because the tree is finite.

Fix a live history \(x\) and a deterministic strategy \(\pi\)
constrained to continue first. Let \(\tau\) count the number of research
choices before terminal action or expiration. Expected resource expense
is \(C_\pi=E_\pi\sum_{s<\tau}c_s\). Delay exposure is
\(T_\pi=E_\pi\tau\geq1\). Let \(R_\pi\) count the action reward only on
histories where action precedes expiration. Then

\[J_b(\pi)-g(x)=R_\pi-C_\pi-g(x)-bT_\pi.\]

The maximum over \(\Pi^+\) is strictly positive if and only if \(b\) is
below the maximum ratio. A negative ratio means that strategy is
inferior even with free additive delay. When \(b\) equals the largest
ratio, stopping ties with a best continuation and is selected by
convention. At histories where continuation is impossible use
\(D=-\infty\). The binary continuation region is nested in \(b\), even
though the optimal experiment need not remain the same. This proof does
not rely on a Gaussian posterior, stationary boundaries, or a one-signal
approximation.

\subsection{B.2 Bellman regret
identity}\label{b.2-bellman-regret-identity}

For a continuation action \(e\) let

\[Q_d(x,e)=-c_e-d+(1-h_e)E_e V_d(x').\]

For a terminal action \(a\), let \(Q_d(x,a)=E[u(a,\theta,x)\mid x]\).
Optimal action payoff \(g\) is the maximum over terminal actions. Extend
\(V\) by zero to the absorbing expired state. At a live continuation
state, the realized period reward is \(-c_e-d\). Its conditional
expectation plus the next-state value equals \(Q_d\). At a terminal
state, realized terminal reward has conditional expectation \(Q_d(x,a)\)
and next value zero.

Subtract the conditional equality from \(V_d(x)\), sum along the
realized finite path, and apply iterated expectations. Intermediate
values cancel. The initial value minus all expected realized rewards is
the expected sum of \(V_d-Q_d\). This yields the formula in separable
benchmark. There is no survival conditioning or exclusion of incomplete
cases. A deadline alters \(\mu\) at every history where its anticipation
changes behavior, and the identity still holds.

\subsection{B.3 State-specific stopping
interventions}\label{b.3-state-specific-stopping-interventions}

Let \(\mu\) be the original analyst policy evaluated at firm costs. At
every selected history \(x\) with \(V_d(x)=g(x)\),
\(J_d(\mu\mid x)\leq g(x)\). Changing the policy to stop there gives at
least its old continuation value. At unselected histories retain the old
action. By induction from terminal histories, improvements in successor
continuation values weakly improve predecessor payoffs, while selected
histories weakly improve because no feasible continuation beats \(g\).
Thus expected firm payoff cannot fall. If some selected reachable
history had a strictly negative old continuation gain and is reached
with positive probability under the new prefix, the improvement is
strict. This argument presumes fixed signal technology and no queue or
effort response.

\subsection{B.4 Worked nonmyopic
example}\label{b.4-worked-nonmyopic-example}

Let current probability of state 1 be .7, signal accuracy .7, no
expiration, correct action reward 100, wrong action zero. One positive
signal has probability .58 and posterior \(.49/.58\); a negative signal
has probability .42 and posterior .5. Expected maximum posterior is
\(.49+.21=.70\). Immediate one-step value of information is zero.

Two further signals yield unconditional state-likelihood pairs:
positive-positive \((.343,.027)\); each mixed path \((.147,.063)\);
negative-negative \((.063,.147)\). Summing maxima gives
\(.343+2(.147)+.147=.784\). Two samples cost two when \(c=1,d=0\), so a
feasible continuation gains 6.4. This is a lower bound on optimal
adaptive continuation value, sufficient to disprove the myopic converse.
The delivered code checks the posterior arithmetic independently of the
Bellman recursion.

\subsection{B.5 Final-period derivation and
boundaries}\label{b.5-final-period-derivation-and-boundaries}

Expected payoff after one experiment followed by action is
\((1-h)(S-L\bar m)-c-d\). Subtract \(S-Lm\). Rearrangement yields
\((1-h)L(m-\bar m)-h(S-Lm)-c-d\). Blackwell precision lowers expected
Bayes classification error at fixed prior and payoff. Hazard lowers the
value through both foregone learning and destruction of immediate
opportunity value. The derivative with respect to \(r\) is nonpositive
even though wrong-action salvage increases the level of immediate
reward.

Higher stakes do not universally favor waiting. For example, at prior
.5, \(r=0,q=.55,h=.2\), the coefficient of \(S\) is \(.8(.55)-.5=-.06\).
With nonnegative costs, continuation is inferior at every positive
stake. Here greater stakes increase the value destroyed by expiration
faster than the value of learning. This is consistent with the reported
high-hazard scenario, but is a separate final-period analytical example.

\subsection{B.6 Delay-cost uncertainty and policy
ranking}\label{b.6-delay-cost-uncertainty-and-policy-ranking}

For fixed policies \(z=0,1\), define \(A_z=R_z-C_z\), so
\(W_z(d)=A_z-dT_z\). For \(d\in[d_L,d_U]\), the contrast interval is

\[[\min\{\Delta A-d_L\Delta T,\Delta A-d_U\Delta T\},
\max\{\Delta A-d_L\Delta T,\Delta A-d_U\Delta T\}].\]

If \(\Delta T=0\), delay price cannot change the ranking. Otherwise the
break-even price is \(\Delta A/\Delta T\), which need not lie in the
admissible interval. This formula values fixed implemented policies. It
is distinct from optimizing \(V_d\) separately for every price.

With heterogeneous \(d_i\in[l_i,u_i]\), use case-level extrema of
\(A_i-d_iT_i\). For a treatment-control difference, subtract the control
upper bound from the treatment lower bound for the lower contrast, and
conversely for the upper. This is conservative when cross-case cost
restrictions are ignored. It does not identify a missing payoff
component merely because its interval is assigned a number.

\subsection{B.7 Optional stopping and
inference}\label{b.7-optional-stopping-and-inference}

The model is a Bayesian economic stopping problem, not repeated
null-hypothesis testing. No signal is a sequence of unadjusted p-values.
The proposed trial uses a fixed final analysis date and fixed enrollment
after feasibility. Exploratory interim monitoring cannot trigger an
ordinary fixed-sample significance claim. If operational needs require
sequential confirmatory monitoring, an alpha-spending plan or an
always-valid method and its assumptions must be chosen and registered
before treatment outcomes are examined. No sequential frequentist method
is implemented in this package.

\subsection{B.8 Computational
verification}\label{b.8-computational-verification}

Python 3 standard-library code implements exact posterior updating on
\((n,k)\), backward induction under perceived costs, and separate
forward evaluation at full costs. In the pre-R1 separable
implementation, 12,002 assertions/checks pass. The parameter-grid count
is 324: three stakes, three recovery fractions, three precisions, four
delay prices, and three hazards. Each compares all 36 nonterminal
histories of an eight-period tree and the initial welfare ordering,
yielding 11,988 checks. Seven frontier comparisons, three posterior
checks, and four derivatives give 12,002 total. The posterior pair
calculation is one of the three posterior checks.

For an independent check, the script enumerates all distinct pairs
\((R-C,T)\) at horizon three, allowing independent policy choices at the
two signal descendants. It retains 15 distinct pairs at the root and
verifies their upper envelope against the Bellman solution. This is
numerical agreement with a finite tree, not a proof for an infinite
horizon. The policy CSV uses a 12-period horizon. Floating-point
comparisons use tolerances; no random seed is required.

\section{Appendix C. Earlier single-case synthetic
illustrations}\label{appendix-c.-earlier-single-case-synthetic-illustrations}

These are retained separable benchmarks, not evidence for the
shared-resource mechanism. The prior is .5, signal accuracy .7, correct
payoff 100, wrong payoff zero, acquisition cost 1, additive delay cost
5, expiration probability .02 per research period, and a 12-sample
finite horizon. The silo internalizes ten percent of additive delay
cost. Values are arbitrary model units, computed by exact backward
induction and separate full-cost policy evaluation.

\begin{center}
\small
\begin{tabular}{@{}p{2.55in}rrrr@{}}
\toprule
Policy or environment & Payoff & Periods & Complete & Accuracy \\
\midrule
Firm optimum & 62.600 & 1.000 & .980 & .700 \\
Silo analyst & 59.002 & 3.294 & .934 & .843 \\
Silo, immediate deadline & 50.000 & 0.000 & 1.000 & .500 \\
Silo, one-sample deadline & 62.600 & 1.000 & .980 & .700 \\
Silo, four-sample deadline & 60.430 & 2.383 & .952 & .785 \\
Firm, minimum two samples & 60.430 & 2.383 & .952 & .785 \\
Firm, recoverable share .8 & 90.000 & 0.000 & 1.000 & .500 \\
Silo, recoverable share .8 & 86.120 & 1.000 & .980 & .700 \\
Firm, correct payoff 500 & 374.065 & 3.294 & .934 & .843 \\
Firm, expiration hazard .20 & 50.000 & 0.000 & 1.000 & .500 \\
\bottomrule
\end{tabular}
\end{center}

Accuracy is conditional on completion. Payoff and completion include
every opportunity. Relative to the silo policy, a one-sample deadline
raises total payoff by 3.598 despite lowering completed-case accuracy.
An immediate deadline lowers payoff to 50. A mandatory second sample can
reduce value by making the intermediate work compulsory while leaving
further continuation endogenous.

The executed legacy script independently checks a three-period
payoff/time frontier against backward induction, verifies
continuation-set nesting on a grid, and checks the nonmyopic
counterexample and final-period derivatives. It reports 12,002 checks.
These results neither identify firm behavior nor establish a novel
single-case index.

The fixed-policy delay-price file separates gross reward less
acquisition expense from active-period exposure. Its break-even
comparison for the frozen one-sample and silo policies is approximately
3.432 delay-cost units. Those bounds hold the full
signal/reward/expiration model and implemented policies fixed.
Reoptimizing at each price answers another question.

The horizon is a substantive finite assumption, not an unverified
infinite-horizon approximation. The R1 main model instead has exactly
two cases, one experiment each, and two slots.

\section{Appendix D. Correlated two-case/two-slot extension: complete
details}\label{appendix-d.-correlated-two-casetwo-slot-extension-complete-details}

\subsection{D.1 Probability space and
filtration}\label{d.1-probability-space-and-filtration}

The latent state vector takes values (0,0),(0,1),(1,0),(1,1) with public
prior mu. Each case has one potential binary signal, conditionally
independent given that vector, with likelihood q\_i for a correct
signal. There are two potential independent survival coins for each
case, independent of the states and signals. A survival coin is used
only when that case is unresolved across that slot. The potential signal
is observed only when its experiment is conducted and its case survives
the corresponding slot. Unused draws integrate out. No latent outcome of
an irrevocable terminal action is revealed during the window;
realization is settled later. This is essential to avoiding unmodeled
free information from action outcomes.

The filtration consists of all observed research assignments, stops,
survival events, and delivered signals. Expiry itself is uninformative
about latent types. For any nonzero-probability observed signal y from
case k,

\[\mu^{k,y}_{ab}=\frac{\mu_{ab}\Pr(y\mid\theta_k)}{\sum_{a',b'}\mu_{a'b'}\Pr(y\mid\theta_k')}.\]

All continuing policies are adapted to this public filtration. At most
two experiments occur in total, never more than one on a particular
case. Case payoffs remain additive even though their conditional
expectations respond to common information.

\subsection{D.2 Final-slot formula}\label{d.2-final-slot-formula}

Fix the current posterior, live set L, and unused-experiment set E.
Stopping every case gives \(G_L=\sum_{i\in L}g_i\). If k is researched,
its own expected resource payoff is \(-c_k-d_k+s_kM_{k|k}\), a gain
b\_k. A passive case i may stop for g\_i or remain unresolved across the
slot. If retained, its expected payoff is

\[-d_i+s_i[(1-s_k)g_i+s_kM_{i|k}].\]

Subtracting g\_i gives

\[-d_i-h_ig_i+s_is_k(M_{i|k}-g_i)=a_{i|k}.\]

This accounts separately for expiry of the recipient and nonarrival of
the source signal. Their hazards are independent, even though the latent
states are correlated. The optimal retention gain is a\^{}+. Summing
this and b\_k, and taking the best single k or stopping all, proves
equation (9). A previously used case may be a recipient and therefore
belongs in L even if absent from E. This is the key difference from the
independent-case enumeration.

If no experiment is selected in the final slot, waiting creates no
signal or action improvement; holding costs and expiry cannot improve
nonnegative expected terminal rewards. Thus final idle policies are
weakly dominated. All researched and retained survivors choose
Bayes-optimal terminal actions after the selected signal, if delivered.
No retrospective reuse of a future signal is allowed for a case that
acted before it.

\subsection{D.3 Initial action exhaustion and rational
solution}\label{d.3-initial-action-exhaustion-and-rational-solution}

There are four retained subsets R. When R is empty there is one active
choice; for each singleton R there are two (no research or that case);
when both are retained there are three. The total is eight. Every case
outside R locks its immediate-action payoff and can neither be reopened
nor researched later. Every case in R pays d\_i for this slot, and a
selected experiment pays c\_k once whether or not its case survives.
Conditional on survivors, if k survived, Bayes-update on its signal; if
not, do not update. These are exactly the terms in (10).

The remaining value is explicitly (9), so (10)-(11) are a finite closed
evaluation of eight rational expressions with maxima, not an unsolved
infinite or general queue recursion. Posterior probabilities and signal
probabilities are rational whenever the primitives are rational. The
supplied solver uses Fraction arithmetic throughout this calculation.

Under independence, the passive mean \(M_{i|k}\) equals g\_i. Thus
passive research-information retention in the last slot is dominated,
including for the previously studied case. The unused case's posterior
is unchanged by the first signal. If scheduled second its optimal
initial retention gain is the independent w\_i, so the eight-action
formula reduces to Theorem 1, with initial idle options dominated as
already proved. This verifies preservation of the original 137.248
independent example; the correlated model does not replace its
assumptions retroactively.

\subsection{D.4 Complete local game}\label{d.4-complete-local-game}

At one remaining slot, define charged own research gains \(b_i-\tau\).
If B has an experiment and \(b_B>\tau\), A's own best release-and-learn
gain is \(a_{A|B}^+\); otherwise it is zero. If A also has an experiment
it claims research precisely when its charged own gain is positive and
weakly exceeds that release gain. A positive tie between early research
and later free learning favors its priority claim; a zero tie favors
immediate action. If A releases or has no experiment, B researches iff
\(b_B>\tau\). Given the active case, each passive owner retains iff its
own a is positive; otherwise it acts before the signal. These are
optimal local actions because passive retention does not alter the
active owner's own signal or resource costs. If no case researches, both
act.

Let the resulting two components be \(L_i^1(\mu,L,E;\tau)\), expressed
at full resource costs without subtracting transfers. For a root choice
(R,k), the actual component is

\[Q_i^L(R,k)=\mathbf1\{i\notin R\}g_i
-\mathbf1\{i\in R\}d_i-\mathbf1\{k=i\}c_i
+E[L_i^1(\mu',L',E_k;\tau)].\]

The private component subtracts the root experiment charge and expected
future experiment charge, each paid only by the active owner. If a
priority toll is applied only when A claims the initial slot, subtract
it only from A's private component on that branch. The transfer is never
subtracted from total resource payoff.

B's follower alternatives after A claims are its immediate action or
retention, giving the set C. If A releases and stays unresolved, B can
act, research now, or wait, giving R. If A acts and releases, B has the
same three own choices with A removed, giving X. These sets of size
2,3,3 exhaust the eight feasible root actions. Choose B's highest
private component in each set, then A's highest private component among
the three resulting routes. This is full backward induction of the
finite perfect-information local game, including off-path choices. No
equilibrium is inferred from a joint-optimal dispatcher policy.

In the full-support synthetic example, at zero charge the three locally
optimal route pairs (A resource payoff, B resource payoff) are:

\begin{itemize}
\tightlist
\item
  A claims, B retains: (78.2, 76.689855).
\item
  A releases and stays, B researches: (90.35532, 64.05).
\item
  A acts and releases, B researches: (60, 64.05).
\end{itemize}

A therefore releases and stays, so the local total is 154.40532. Its
claim route has total 154.889855 but is privately worse for A. Even that
route does not achieve the efficient 159.69508928 because B declines a
later experiment with negative own gain. Adding the R1 marginal-only
priority toll .6405 to the claim route does not change the equilibrium.
This computation exposes both initial free-riding and a supported
continuation underresearch boundary.

\subsection{D.5 Signed continuation incentive and full payoff
alignment}\label{d.5-signed-continuation-incentive-and-full-payoff-alignment}

At a last-slot history with only i's experiment available, the joint
gain is b\_i+a\_j\textbar i\^{}+. At a nonnegative fee the active owner
will not undertake an experiment with nonpositive own b\_i. If the sum
is positive despite that, no such fee implements. A conditional credit
equal to a\_j\textbar i\^{}+ makes the owner's gain equal to the joint
comparison. A strict credit above -b\_i is required for strict research
if b\_i\textless0, with equality permitting the prescribed action tie.
This proves Theorem 6's exact boundary.

This final-slot credit is not enough to certify the root equilibrium.
A's release-and-free-learn payoff and B's retention responses can also
change. The full finite-game solution above must be recomputed under any
proposed transfer scheme, including future credits. The paper does not
claim that the last-slot credit alone attains the optimal root value.

For the stronger theoretical team-payoff contract, let realized resource
payoffs, including all own costs and terminal reward, be X\_A and X\_B.
Give i a signed additional internal transfer X\_j-K\_i, where K\_i is an
initial fixed constant independent of subsequent choices and signals.
Its private total becomes X\_A+X\_B-K\_i. At every public decision,
maximizing conditional expected private payoff is maximizing the same
total resource value. The finite game's backward induction then selects
the joint-optimal feasible continuation at every history, and hence
(9)-(11). All eight root choices and all final retention plans are
attainable through the defined sequential moves, so no central-feasible
action is excluded by the priority right when both owners share this
objective.

This is an elementary sufficient alignment result, not a new
private-type mechanism. Transfers may be positive or negative, need not
balance within owners, and require observable eventual payoffs and
authority over internal marginal incentives. Adding both transfers to
social payoff would be wrong: the firm objective is still X\_A+X\_B, not
twice that quantity. A fixed constant changes distribution, not
decisions. These assumptions are stronger than those of the restricted
one-time priority fee, and the construction is not asserted to satisfy
the published dynamic-pivot participation or truthfulness conditions
under correlated private information.

\subsection{D.6 Exact example and independent
enumeration}\label{d.6-exact-example-and-independent-enumeration}

\begin{sloppypar}
For mu=(.58,.02,.02,.38), q\_A=.8, q\_B=.95, c\_A=1, c\_B=30,
h\_A=h\_B=.01, S\_A=S\_B=100, and zero d and r, all state and relevant
signal branches have positive probability. Conditional on A surviving,
its signal is 1 with probability 11/25 and 0 with probability 14/25.
After signal 1 the posterior is (29/110,1/110,2/55,38/55). After signal
0 it is (29/35,1/35,1/140,19/140).
\end{sloppypar}

If both survive and the first signal is 1, the final research plan keeps
both unresolved and researches B. At this posterior, g\_A=800/11,
g\_B=70. B's own gain is -119/20, while A's cross-information gain is
9401/550. The total increment is 12257/1100. Expected final case payoffs
are 4491/50 for A and 1281/20 for B, totaling 15387/100. With a first
signal 0, both act for a total 1185/7. These values include the
final-slot costs and hazards.

The six first-slot branches have probabilities 1/10000 (neither alive),
99/10000 (only B alive, no delivered A signal), 693/125000 (only A
alive, signal 0), 1089/250000 (only A alive, signal 1), 68607/125000
(both alive, signal 0), and 107811/250000 (both alive, signal 1). They
sum to one. Weighting the final values and subtracting first-slot
acquisition/holding costs gives 249523577/1562500 exactly. The root
comparison with every other retained-set/active choice is supplied in
JSON and the main table.

The independent enumerator pre-draws physical latent states, both
potential signals, and both survival coins per case. It evaluates actual
terminal correctness on each path, using a separately calculated
likelihood-based action at the information available when that action is
locked. The evaluator does not call the rational final-slot formula or
root solver. It enumerates all root actions and all observable-history
contingent final plans, including retention of the already researched
case. Final idle plans are excluded only by the analytic dominance
proof; root idle plans are enumerated.

For positive-probability interior branches the number of final
schedule/retention maps per root action, in order empty, idle A, active
A alone, idle B, active B alone, idle both, active A with both, active B
with both, is 1,2,1,2,1,20,18,18, totaling 63. For idle both, both-live
states allow five plans (stop; A active with B stopped or retained; B
active with A stopped or retained), each singleton survivor allows two,
giving 5\emph{2}2=20. For active A with both retained, each both-live
signal node allows three plans and the B-only survivor node allows two,
giving 3\emph{3}2=18. Other nodes have no remaining experiment or no
case. Branches of zero probability reduce counts. Bayes-optimal terminal
actions are imposed, not counted as distinct deliberately mistaken
policies.

There are exactly three named checks, with 63,30,63 maps, totaling 156
output rows. They are the full-support correlated example, a
perfect-correlation zero-expiry boundary, and the product prior with the
same marginal probabilities and other primitives. Both all-root
comparisons and implementation boundary checks are executed, 37 in
total. This is a finite exact calculation and enumeration, not a Monte
Carlo simulation or a larger queue study.

\section{Appendix E. Two-credit implementation: exact regional
proof}\label{appendix-e.-two-credit-implementation-exact-regional-proof}

\subsection{E.1 Fixed institution, no omitted
payments}\label{e.1-fixed-institution-no-omitted-payments}

This appendix proves Theorems 7-8 on I={[}17/100,23/100{]}. It retains
Appendix D's latent binary states, public signal histories, independent
.01 expiration, one experiment per case, local priority rights, and
immediate-action options. C is earned by conducting initial A research,
not by retaining both or by successful survival. B's credit requires
that A was initially researched, both survived the first slot, A's
delivered signal was 1, and B conducts its final experiment. The payment
is earned before its second survival coin. No credits arise after
initial B research, initial idling, or terminal action. Credits are sunk
once earned and do not change subsequent terminal-action preferences. No
experiment can be repeated to collect another credit.

The zero-own-gain research/action tie selects action. A positive initial
claim/release tie selects its priority claim, consistently with the
original local institution. These are different ties. Conditional
Bayes-action ties can use action 1 without affecting expected payoffs.
Weak implementation means the outcome chosen by these conventions;
strict supported implementation means strict root and discretionary
on-path research/retention incentives. It does not require every unused
off-path option to have a strict ranking.

\subsection{E.2 Joint policy and complete exact interval
certificate}\label{e.2-joint-policy-and-complete-exact-interval-certificate}

Write s=99/100 and h=1/100. The target policy has the fixed
research/retention decisions in Section 8.4. Terminal actions are also
fixed on I: if only A survives its first experiment, act according to
that signal; after both survive signal 0, both choose action 0; after
both survive signal 1 and B is studied, both surviving cases choose B's
signal, except that A uses its initial signal if B expires and no second
signal arrives. If A expires initially, B's one experiment uses the
prior and its own signal. All these terminal choices are Bayes optimal
throughout I. This follows by comparing the unnormalized masses of each
binary state after each observed signal. The only two-signal
Bayes-action switch inside I is at 41/225 for A after (Y\_A,Y\_B)=(0,1),
a signal pair the target policy never generates because it stops after
Y\_A=0.

The target's fixed physical resource components can be derived without
posterior division. The relevant joint probabilities are

\[\Pr(\theta_A=0,Y_A=0)=.48,\quad
\Pr(\theta_A=1,Y_A=1)=.32,\]
\[\Pr(Y_A=1,Y_B=\theta_A)=.202+.9\delta,\quad
\Pr(\theta_B=0,Y_A=0)=.336+.6\delta.\]

Consequently

\[X_A=-1+100\{sh(.8)+s^2(.48)+s^3h(.32)
+s^4(.202+.9\delta)\},\] \[X_B=sh(64.05)+100s^2(.336+.6\delta)
+s^2(.44)(64.05). \tag{E1}\]

The A cost is paid once even if it expires. B's research payoff 64.05
already includes its final acquisition cost 30 and survival-weighted
reward .99(95). The B-only branch uses that value, while the
both-surviving signal-1 branch has probability s\^{}2(.44). Thus neither
cost nor lost opportunity is counted twice. Expanding (E1) gives (14),
so the target value is affine in delta.

For completeness, the maximum physical payoff at every root choice is
certified exactly at both interval endpoints:

\begin{longtable}[]{@{}lrr@{}}
\toprule\noalign{}
Retained set, initial assignment & delta=17/100 & delta=23/100 \\
\midrule\noalign{}
\endhead
\bottomrule\noalign{}
\endlastfoot
Neither, none & 120 & 120 \\
A only, idle & 68709/500 & 68709/500 \\
A only, A & 691/5 & 691/5 \\
B only, idle & 246819/2000 & 246819/2000 \\
B only, B & 2481/20 & 2481/20 \\
Both, idle & 36853237179/250000000 & 38582309997/250000000 \\
Both, A & 30486421447/200000000 & 161147685689/1000000000 \\
Both, B & 729696081/5000000 & 312339/2000 \\
\end{longtable}

These fractions are finite sums over latent states, potential binary
signals and the four potential independent survival coins. At each root
the evaluator enumerates every final schedule/retention map, with
terminal Bayes actions evaluated separately from the joint solver. There
are 63 maps per prior. At both endpoints and the witness, the target is
the unique maximizing map; the closest distinct map is lower by
8019/200000=.040095. This is the loss from omitting the worthwhile
B-only continuation on its positive-probability survival branch.

To prove the continuum, fix any feasible pure physical policy, including
its terminal actions. Its expected payoff is affine in delta because the
joint prior is affine and signal/survival likelihoods are fixed. The
maximum over terminal actions or any collection of policies is therefore
convex. For any competing root choice, its optimized payoff is convex,
while the target payoff (E1) is affine. The target strictly exceeds it
at both endpoints by the table, hence strictly exceeds it throughout the
closed interval by the defining convexity inequality. The same reasoning
compares every distinct final schedule/retention map under the target
root: its terminal-action-optimized value is convex, and the target
exceeds each at both endpoints. This proves strict optimality of the
target map on I. No extrapolation from a sampled grid or unproved
monotonicity is used. Appendix D's final-idle dominance removes only
decisions that cannot improve a payoff; root idle choices are explicitly
included.

\subsection{E.3 All final-slot local
continuations}\label{e.3-all-final-slot-local-continuations}

After initial A research and survival of both cases:

\begin{enumerate}
\def\labelenumi{\arabic{enumi}.}
\tightlist
\item
  If Y\_A=1, B's own uncredited research gain is -K. Its gain including
  the observable-history credit is kappa-K. The zero-gain tie selects
  action, so research occurs exactly when kappa\textgreater K. A's
  passive cross-information gain is (4410450 delta-594259)/22000,
  strictly positive on I, and it strictly retains if research occurs.
  Otherwise both act. Joint research is strictly valuable because the
  sum is (1410450 delta-65159)/22000\textgreater0.
\item
  If Y\_A=0, B's own gain is -3(5000 delta-189)/140\textless0 and no
  credit is eligible. Thus it acts, and A acts as well. A's hypothetical
  benefit from a B experiment cannot itself make B conduct it.
\end{enumerate}

If A expires but B survives, expiry is uninformative and B researches at
gain 81/20. If only A survives, no experiment remains and it acts using
its delivered signal. If neither survives, there is no choice or further
payoff. All six positive-probability first-slot histories are thereby
covered. Payment C was already earned and is irrelevant to all these
continuations.

Off the target route, after initial B research, no credit is ever
eligible. If both survive B's signal 0, A's own research gain is -2(225
delta+236)/295\textless0, and both act. After B's signal 1, A researches
exactly when delta\textless delta\_0. On this profitable branch its own
gain is (7831-45000 delta)/205, and B acts before the experiment because
its passive gain is -38/41. At delta\_0 A's zero gain selects action.
For delta\textgreater delta\_0 research remains unprofitable; beyond
41/225 the uncredited own gain has the expression -(450 delta+287)/205,
also strictly negative. B's passive cross gain after its signal 0 is
-57/59. If B expires and A survives, A researches with uncredited gain
91/5. If only B survives its own first experiment, it acts. These are
all B-first survival/signal continuations.

After initial idle retention of both cases, both experiments remain. If
both survive, A's own immediate research gain is 91/5, whereas releasing
to B gives A the larger cross gain

\[a_{A|B}(\mu)=\frac{99^2}{100^2}100[(.518+1.8\delta)-.6]-.6>91/5.\]

The inequality holds at 17/100 and increases thereafter. B's own gain is
81/20\textgreater0, so A releases, B researches and A retains. The
own/passive components are precisely those of the local game, not the
joint-optimal final assignment. If only A or B survives, that case
conducts its own profitable experiment. With a single case retained
initially, the same singleton continuation applies; a case researched
initially cannot repeat research. With no live cases, nothing occurs.
This exhausts initial-idle, singleton and empty-root continuations as
well as both research roots.

For transparency about unused ties: the Bayes-action tie at 41/225 is on
an unchosen two-signal route; A's hypothetical retention tie after
initial Y\_A=0 and a subsequently conducted B experiment is at
delta=137947/735075, because the cross gain above 41/225 is 3(735075
delta-137947)/14000. B never chooses that uncredited experiment on I.
Neither tie changes (17) or funding. A stronger requirement that every
conceivable off-path action and terminal Bayes choice be strict would
fail at these points and at delta\_0 without further instruments. The
stated strict supported implementation deliberately does not claim that
stronger property.

\subsection{E.4 Root follower choices, including failed-continuation
credits}\label{e.4-root-follower-choices-including-failed-continuation-credits}

If A initially claims, B can act for 60 or retain. Let

\[Y_B^0=50.815215+117.612\delta.\]

This is its resource/private payoff from retention if it does not
conduct research on the both-live signal-1 branch. It includes the
uncredited B-only experiment after A expires. It is strictly above 60
throughout I. If kappa\textgreater K, its retained private payoff is
\(X_B+p\kappa=Y_B^0+p(\kappa-K)>60\). If kappa\textless=K, no eligible
experiment occurs and the retained payoff is Y\_B\^{}0\textgreater60.
Thus B strictly retains for \textbf{every} pair of nonnegative
instruments, not only after imposing the target policy. This also proves
the retention slack at the witness's limiting credit.

If A releases but stays unresolved, B may (i) act immediately, leaving
only A to wait, (ii) conduct its experiment with both retained, or (iii)
wait with A. Its respective own payoffs are

\[60,\qquad64.05,\qquad .99(64.05)=63.4095.\]

The third value follows from the actual no-signal final local game just
derived, including the cases in which A expires first. It is not
inferred from a centrally imposed continuation. B strictly chooses (ii).
If A acts and releases, B has exactly the same own three values and
again researches initially. These comparisons fully resolve the follower
sets C,R,X of Section 8.2, including idle and immediate-action
alternatives.

A's private payoffs for claim, release-and-retain, and act-and-release
are therefore

\[\begin{cases}(X_A+C,R_A,60),&\kappa>K,\\
(78.2+C,R_A,60),&0\leq\kappa\leq K.
\end{cases} \tag{E2}\]

The second branch exists only when K is nonnegative. Without the later B
experiment, A's research payoff is .99(80)-1=78.2. In both branches B's
root retention is optimal. Since R\_A\textgreater60 and
R\_A\textgreater78.2 on I, the actual local equilibrium is A-first
exactly when C is at least R\_A minus its corresponding research
component in (E2); otherwise it releases and waits. A's act-and-release
option is feasible but strictly dominated by release-and-retain here.
This solves root behavior also for instrument pairs that fail
implementation.

To derive R\_A, if B expires initially and A survives, A researches for
78.2. If both survive, without A's further experiment it acts according
to B's signal, whose unconditional correctness probability is .518+1.8
delta. After B's signal 1, its possible own profitable research adds the
positive part derived above. Hence

\[R_A=sh(78.2)+100s^2(.518+1.8\delta)
+s^2(.41)[(7831-45000\delta)/205]^+,\]

which is (18). Subtracting (E1) yields the two pieces (16). The low
piece has positive intercept and slope; the high piece is increasing and
agrees with the low piece at delta\_0. F is positive everywhere on I.
Combining (E2) with B's strict research condition proves both directions
of (17), not merely a sufficient deviation check. Once the continuation
is induced, kappa has no further effect on A's resource component, so
even an arbitrarily large continuation credit cannot replace C. A
nonnegative initial toll makes its strict free-riding preference
stronger.

\subsection{E.5 Funding constraints and endpoint
interpretation}\label{e.5-funding-constraints-and-endpoint-interpretation}

The payment random variable on the implemented policy is
\(C+\kappa\mathbf1_{\mathcal H}\), where
\(\Pr(\mathcal H)=p=s^2(.44)>0\). The event is research after the first
signal-1/both-survive history; it does not require either second-slot
survival. Both events of receiving C alone and of receiving C+kappa have
positive probability. Thus (19) is exact, including the payment after
initial A expiry. No eventual realized reward need be observed for
settlement.

Put \(k=K^+\). The implementation set is the product of C\textgreater=F
and the half-line kappa\textgreater=0,kappa\textgreater K; for strict
supported implementation replace the first inequality by
C\textgreater F. Consequently the weak expected/pathwise infimum
functions are piecewise F below delta\_1 and F+pK/F+K at or above
delta\_1, with F itself changing slope at delta\_0. At delta\_0 the two
F formulas coincide; kappa=0 is feasible because K\textless0. At
delta\_1, kappa=0 fails by the zero-gain action tie, so the continuous
infimum ceases to be attained. At both closed endpoints of I the
respective formulas and tie conventions apply without taking an unstated
open neighborhood. No claim is made just outside I or at independence,
where the target pattern need not persist.

For weak implementation, a single cap equal to its infimum is feasible
exactly when delta\textless delta\_1. For strict supported
implementation no equality cap is feasible anywhere on I. Under
simultaneous expected/pathwise caps, both corresponding slack or weak
inequalities are required. Sufficiency follows by choosing sufficiently
small increments in each payment, rather than treating an expected limit
as permission to exceed a pathwise limit. Since both outlays are
increasing in C and kappa, there is no unaccounted trade-off permitting
a smaller initial payment in exchange for a larger continuation payment
after the efficient branch is induced. These are gross funding
requirements for this specified class, not firm resource costs or
minimum funding over all possible mechanisms.

\section{P.1 Status, unit, and estimands before
methods}\label{p.1-status-unit-and-estimands-before-methods}

No partner, data agreement, approved treatment, recruitment,
observations, or experiments exist. This is a prospective protocol. The
theoretical model has two one-experiment cases and two slots, with
independent and correlated binary-state specifications. The study tests
its shared-capacity mechanism in eligible operational queues, not its
literal numerical calibration.

Primary unit: an analyst-queue episode with a frozen intake cohort and a
defined capacity window. A queue consists of all cases sharing the same
analyst/reviewer resource. An episode can contain more than two cases;
theoretical optimality is not asserted for that larger queue. A matched
two-case substudy is possible only if the institution can document a
bounded two-slot decision window and reserve capacity without displacing
uncounted cases. Cases excluded from the cohort but affected by reserved
capacity must be followed as spillovers.

Primary estimand: intention-to-treat effect of offering an
opportunity-loss scheduling aid on total 90-day net payoff summed over
the episode's registered cases, averaged equally across a crossed
evidence-provision assignment. The primary target is the mean episode
total at the prespecified cohort definition, not a mean among completed
purchases. Report case-weighted results secondarily. Unequal episode
size and any weighting are frozen before assignment.

If defensible monetary bounds are incomplete, the primary estimand is
interval valued. Declare a favorable welfare effect only when its lower
confidence bound is positive across the prespecified admissible
valuation set. A faster queue with worse total expected payoff is not a
successful treatment.

Secondary confirmatory family, with Holm adjustment across seven
contrasts: evidence provision effect on total payoff; aid-by-evidence
interaction on total payoff; aid effect on all-worker analysis minutes;
aid effect on summed unresolved case-hours through day 90; aid effect on
cases completed by day 30; aid effect on expired cases; aid effect on
total rework and switching expenditure through day 90. Case reversal
rates and mechanism interactions are exploratory and explicitly labeled.

\section{P.2 Feasibility and
eligibility}\label{p.2-feasibility-and-eligibility}

Seek a partner with routine, non-safety-critical procurement decisions,
reliable intake for purchases and rejections, logged evidence requests,
a shared analyst roster, verified quote expiry, and linked operational
outcomes. Exclude safety-critical, legally mandated, personnel,
emergency, and strategic high-value reviews. The partner defines
exposure caps, authority limits, and reversible-action terms before
enrollment. No such caps are invented here.

A read-only feasibility audit must establish: the actual shared-resource
boundary; how managers currently claim priority; whether immediate
action really avoids analyst use; packet production durations; existing
mandatory reviews; rejected/unresolved intake coverage; capacity window
commitments; and spillovers to adjacent queues. If terminal approval
itself uses the same scarce reviewer, the two-slot research-only model
does not fit without redesign. Record that mismatch rather than silently
treating both services as the same.

Public BPI Challenge 2019 organizer documentation verifies a
procurement-process event-log candidate. Repository/file access and
license were not verified; no schema or rows were analyzed. It lacks a
verified research-queue mapping in this work. It may inform process
measurement if later accessed under its actual terms, but cannot
identify the treatment effect without intake, resource links, and
economic outcomes.

\section{P.3 Design and treatments}\label{p.3-design-and-treatments}

Use a parallel 2-by-2 design at the whole independent analyst-pool
level. If too few pools exist, consider randomly assigned pool-period
blocks only after demonstrating feasible washout and tracking all
carryover cases; this changes assumptions and requires a new
registration. Do not randomize individual cases independently in one
shared queue and then assume no interference.

Evidence factor: provide one frozen incremental packet for each eligible
case, versus existing information. Packet production time, source costs,
and worker effort count in the resource ledger. Keep purchasing
alternatives, legal information, price opportunities, and formal
authority comparable. Display design is held constant conditional on
provided facts. A compressed-display substudy, if later chosen, must
preserve accessible content and has an attention estimand rather than an
acquisition estimand.

Scheduling factor: offer a queue-level opportunity-loss aid versus
existing priority practice. At a review point, the aid records for each
case: current best action; available incremental experiment; credible
expected action improvement; acquisition time; option to act
immediately; consequences of waiting; quote expiry; recoverability; and
the identity of other cases displaced by the next analyst slot. It
distinguishes a research value estimate from a deferral-loss estimate.
It displays a proposed research order and explicit act-now or
retain-for-later decisions, with documented overrides.

Record prespecified cross-case prior links, signal-source identities,
posterior revisions in both cases, and whether an already researched
case remains unresolved to receive another case's signal. Do not infer
independence merely from separate case records. Correlated episodes
require a joint prior and signal model; applying the independent u-score
to updated marginals is not justified. The final-slot information credit
is a theoretical incentive calculation, not an authorized transfer
treatment.

In verified independent two-case one-packet episodes, the aid can
calculate b, w, and u under transparent parameter inputs. In larger
queues it is only a documented scheduling recommendation, not a
theorem-backed globally optimal index. Uncertain inputs must produce
sensitivity ranges, not a falsely precise rank. Record when ranks are
indeterminate over the full uncertainty set.

The offered aid does not implement real marginal incentives
automatically. No effect is labeled a toll-mechanism treatment effect.
If management later authorizes actual internal charges, register a
distinct arm specifying who pays, when, whether later service is
uncharged, whether the charge changes the marginal budget, and whether
the analyst's priority right is enforceable. A generic charge per minute
is not the posted first-slot charge in Theorem 3. No actual charge,
binding deadline, or authorization change is part of the present study.

\section{P.4 Interference and
identification}\label{p.4-interference-and-identification}

Treatment directly changes delay on other cases; that spillover is the
mechanism. The primary total episode payoff includes it. Randomization
identifies the policy-offer effect if cohort registration is unaffected
by assignment, outcomes are measured comparably, and cross-pool
spillovers or carryover are either negligible or included in a
prespecified exposure design.

Use staffing and resource rosters to identify shared analysts,
reviewers, vendors, stock constraints, and work redirected outside the
pool. Log both treated and displaced cases, including unserved requests.
Estimate exploratory effects by exposure to treated competing cases, but
do not call those causal direct effects without a registered
interference model. If resource sharing is organization-wide, the
independent-pool design may be infeasible. The correct response is
redesign, not inflated apparent sample size.

Analyze with a saturated factorial regression or assignment-based
contrast at episode level, retaining randomization strata and frozen
baseline load, case mix, and holding exposure covariates. If episodes
repeat within pools, inference clusters at the assignment pool. Use
randomization inference consistent with the actual assignment, with
small-cluster robust inference as a companion. Define the aid main
contrast as equally averaged across the two evidence levels.

Do not control in the main regression for post-treatment completed-case
status, evidence requested/opened, research minutes, priority rank, or
elapsed time. These are mediators/outcomes. Pre-treatment difficulty
measures include novelty, baseline available evidence, stakes, number of
alternatives, resource requirement, approval layers, and time
sensitivity.

The mechanism prediction is that changing priority rights or accounting
for a rival's deferral loss matters most when local and joint orderings
disagree. Define disagreement using frozen pre-treatment inputs and a
documented valuation uncertainty set. Larger effects in queues with
split budgets or severe load are descriptive heterogeneity unless those
features are randomly manipulated. A future randomized display of
verified competing-case costs could sharpen a mechanism test while
holding the focal packet fixed.

\section{P.5 Data dictionary and
timing}\label{p.5-data-dictionary-and-timing}

\begin{longtable}[]{@{}
  >{\raggedright\arraybackslash}p{(\columnwidth - 6\tabcolsep) * \real{0.2500}}
  >{\raggedright\arraybackslash}p{(\columnwidth - 6\tabcolsep) * \real{0.2500}}
  >{\raggedright\arraybackslash}p{(\columnwidth - 6\tabcolsep) * \real{0.2500}}
  >{\raggedright\arraybackslash}p{(\columnwidth - 6\tabcolsep) * \real{0.2500}}@{}}
\toprule\noalign{}
\begin{minipage}[b]{\linewidth}\raggedright
Variable
\end{minipage} & \begin{minipage}[b]{\linewidth}\raggedright
Definition and units
\end{minipage} & \begin{minipage}[b]{\linewidth}\raggedright
Source/timing
\end{minipage} & \begin{minipage}[b]{\linewidth}\raggedright
Main concern
\end{minipage} \\
\midrule\noalign{}
\endhead
\bottomrule\noalign{}
\endlastfoot
pool\_id & All cases sharing analyst capacity & Baseline roster & Hidden
shared workers \\
episode\_id & Frozen cohort and capacity window & Intake before
assignment & Selective episode construction \\
case\_id & Unique decision, bundled lines resolved & Intake &
Split/duplicate cases \\
eligible & Pre-treatment inclusion rule & Intake attributes & Strategic
exclusions \\
assigned\_E, assigned\_S & Evidence/aid assignments & Randomization
ledger & Offer versus uptake \\
competing\_ids & Cases displaced by a requested slot & Resource ledger &
Unrecorded opportunity loss \\
priority\_owner & Existing first-slot right and override authority &
Baseline process & Informal authority \\
analyst\_available & Research capacity by slot and worker & Staffing log
& Capacity altered by treatment \\
request\_time & First research/decision intake & System log &
Backdating \\
research\_start\_end & Work on each experiment & Work log/time sample &
Offline work \\
passive\_exposure & Unresolved waiting time without own research &
Derived log & Calendar time versus active effort \\
action\_time\_type & Purchase, reject, expire, or unresolved &
Authorization/terms & Missing rejections \\
expiry\_time\_terms & Verified quote/opportunity expiry and extension &
Vendor terms & Endogenous extension \\
prior\_risk & Frozen uncertainty assessment & Pre-treatment rubric &
Unvalidated probabilities \\
packet\_quality & Likelihood/precision evidence where estimable &
External historical validation & Availability not use \\
stakes\_recovery & Documented action loss and recoverable share/cost &
Baseline contracts & Recovery treated as free \\
all\_worker\_minutes & Focal and displaced workers' active time & Linked
work sampling & Unobserved displacement \\
terminal\_net\_value & Realized action margin before other exclusive
costs & Finance/operations & Missing unchosen outcome \\
rework\_switching & Downstream corrections, reversal and replacement
cost & 90-day follow-up & Delayed posting \\
residual\_holding & Cost schedule not counted in margins/expiration &
Independent costing & Double counting \\
cost\_bounds & Justified lower/upper costs and dependencies & Before
outcomes & Favorable valuation selection \\
unresolved\_90 & Registered case still open at day 90 & Follow-up &
Completer selection \\
episode\_bounds & Sum of all case and resource payoff bounds & Derived
ledger & Omitted cases and external work \\
\end{longtable}

This is a proposed schema, not verified columns from a firm dataset.
Independent validation is needed before treating stated confidence as a
calibrated prior or a packet's quality as a likelihood. Data
availability, precision, actual evidence use, capacity rights, and
outcomes are distinct variables.

\section{P.6 Welfare ledger, missing outcomes, and bounded
costs}\label{p.6-welfare-ledger-missing-outcomes-and-bounded-costs}

Each resource cost and consequence has one ledger category. Acquisition
expense counts actual production and analyst work once. Holding exposure
counts residual lost capacity or operating costs not already counted in
action margin or expiration losses. Expired cases have no realized
opportunity payoff under the literal model, but real organizations may
retain salvage; the field ledger must use verified institutional terms
instead. Transfers for hypothetical internal priority charges are
excluded from organization-wide payoff.

The theoretical externality u\_j is not added to observed downstream
loss. It is an opportunity-cost comparison used to schedule; realized
policy payoff already includes the rival case. Adding both would double
count. Likewise, a nominal time saving is not necessarily cash savings;
report physical minutes and scarce-capacity shadow-price bounds
separately.

At a fixed implemented policy, mean payoff is evaluated for all enrolled
cases over the same horizon. Random assignment identifies total realized
policy value without observing every unchosen supplier's quality. It
does not identify decision-specific regret relative to an oracle. Common
independent ex post market information may support a secondary quality
measure only under a validated mapping to economic outcomes.

Report unresolved and missing-payoff fractions by arm. Never compare
mean delay only among completed cases. Sum restricted case-hours through
the fixed horizon, with explicit rejection treated as resolution and
expiration as a competing event. If missing case payoff is
institutionally bounded by {[}L\_i,U\_i{]}, use those bounds. Without
defensible limits, do not assign zero or simulate informative losses.
Report worst-case mean contrasts and sensitivity to missingness. MAR
imputation may be secondary only with assumptions stated.

For fixed policy z and known non-delay primitives, uncertain residual
costs give \(W_z=A_z-\sum_i d_iT_{iz}\). Endpoint bounds hold that
policy and the full signal/reward/expiration model fixed. They do not
establish robustness to model error. With uncertain priors, recovery,
precision, or hazards, define a full joint uncertainty set with
documented dependencies and recompute every candidate policy contrast.
Require uniform nonnegative lower comparison for a certified ordering;
otherwise report ambiguity. No such real-data uncertainty set is
currently estimated.

\section{P.7 Precision and feasibility
planning}\label{p.7-precision-and-feasibility-planning}

Illustrative arithmetic only: with 120 genuinely independent analyst
pools, 20 episodes per pool, balanced four-cell assignment, and
episode-payoff ICC .05, the design effect is 1+19(.05)=1.95. The
normal-approximation 80\% power, two-sided 5\% aid-main-effect MDE is
\((1.96+.84)\sqrt{4(1.95)/2400}=.160\) episode-payoff standard
deviations. The corresponding factorial interaction MDE is .319 SD. At
ICC 0 and .10, main-effect MDEs are .114 and .195 SD. These are
hypothetical pools and episodes, not observed sample sizes, variances,
or power estimates.

The independent assignment unit must be real. Twenty cases in one
episode are not twenty independent queue experiments, and repeated
episodes with the same analyst require clustered inference. Unequal
episodes, fewer pools, carryover, uptake failures, and wide payoff
bounds worsen precision. Final power requires a partner's baseline
outcome variance, ICC, assignment structure, and economically meaningful
currency or capacity effect, followed by simulation under the actual
design. None is available.

A feasibility pilot measures only nuisance quantities and operational
limits, not a treatment-effect-driven outcome choice. Freeze final
enrollment, cohort rules, hypotheses, valuation bounds, and analysis
dates before examining treatment outcomes. No unadjusted repeated
p-value peeking is planned. Any confirmatory interim analysis requires a
separately registered alpha-spending or always-valid method.

Go/no-go gates: complete intake; real independent capacity pools;
auditable priority rights; comparable alternatives and capacity; bounded
stakes; linked all-case outcomes; displaced-work coverage; and
institutional authorization. Failure triggers redesign rather than a
nominal larger N.

\section{P.8 Unpopulated reporting
template}\label{p.8-unpopulated-reporting-template}

\begin{longtable}[]{@{}
  >{\raggedright\arraybackslash}p{(\columnwidth - 10\tabcolsep) * \real{0.1667}}
  >{\raggedright\arraybackslash}p{(\columnwidth - 10\tabcolsep) * \real{0.1667}}
  >{\raggedright\arraybackslash}p{(\columnwidth - 10\tabcolsep) * \real{0.1667}}
  >{\raggedright\arraybackslash}p{(\columnwidth - 10\tabcolsep) * \real{0.1667}}
  >{\raggedright\arraybackslash}p{(\columnwidth - 10\tabcolsep) * \real{0.1667}}
  >{\raggedright\arraybackslash}p{(\columnwidth - 10\tabcolsep) * \real{0.1667}}@{}}
\toprule\noalign{}
\begin{minipage}[b]{\linewidth}\raggedright
Assignment
\end{minipage} & \begin{minipage}[b]{\linewidth}\raggedright
Registered episodes
\end{minipage} & \begin{minipage}[b]{\linewidth}\raggedright
All cases
\end{minipage} & \begin{minipage}[b]{\linewidth}\raggedright
Unresolved fraction
\end{minipage} & \begin{minipage}[b]{\linewidth}\raggedright
Total minutes
\end{minipage} & \begin{minipage}[b]{\linewidth}\raggedright
Payoff lower/upper
\end{minipage} \\
\midrule\noalign{}
\endhead
\bottomrule\noalign{}
\endlastfoot
Standard packet, local order & Not observed & Not observed & Not
observed & Not observed & Not observed \\
Added packet, local order & Not observed & Not observed & Not observed &
Not observed & Not observed \\
Standard packet, queue aid & Not observed & Not observed & Not observed
& Not observed & Not observed \\
Added packet, queue aid & Not observed & Not observed & Not observed &
Not observed & Not observed \\
\end{longtable}

No empirical estimates, assignments, p-values, or implemented
organizational outcomes are supplied. Synthetic model results belong to
the theory section only.

The two-credit result is theoretical only. A future feasibility audit
could distinguish expected internal credit authorization from gross
pathwise funding capacity and verify that research and the specified
public signal history are observable. It supplies no authorization to
pay credits or add an experimental treatment, and no real transfer data
exist.

\end{document}
